\documentclass[11pt,a4paper]{article}
\usepackage[T1]{fontenc}
\usepackage[utf8]{inputenc}
\usepackage{lmodern}
\usepackage[margin=27mm,headheight=14pt]{geometry}
\usepackage{amsmath,amssymb,amsthm,mathtools}
\usepackage{booktabs,longtable,array,enumitem,microtype}
\usepackage[dvipsnames]{xcolor}
\usepackage{listings,fancyhdr}
\usepackage{xurl}
\usepackage[hypertexnames=false,colorlinks=true,linkcolor=MidnightBlue,citecolor=MidnightBlue,urlcolor=MidnightBlue]{hyperref}
\usepackage[nameinlink,noabbrev]{cleveref}
\definecolor{ink}{HTML}{18364C}
\definecolor{pale}{HTML}{F3F6F8}
\newcommand{\N}{\mathbb N}
\newcommand{\Z}{\mathbb Z}
\newcommand{\Q}{\mathbb Q}
\newcommand{\R}{\mathbb R}
\newcommand{\ipra}{\mathrm{iPRA}}
\newcommand{\PR}{\mathrm{PR}}
\newcommand{\eps}{\varepsilon}
\newcommand{\one}{\mathbf 1}
\newcommand{\norm}[1]{\left\lVert#1\right\rVert}
\newcommand{\abs}[1]{\left|#1\right|}
\newcommand{\dw}{d_w}
\newcommand{\gap}{\gamma}
\newtheorem{theorem}{Theorem}[section]
\newtheorem{lemma}[theorem]{Lemma}
\newtheorem{proposition}[theorem]{Proposition}
\newtheorem{corollary}[theorem]{Corollary}
\theoremstyle{definition}
\newtheorem{definition}[theorem]{Definition}
\newtheorem{example}[theorem]{Example}
\theoremstyle{remark}
\newtheorem{remark}[theorem]{Remark}
\newtheorem{warning}[theorem]{Boundary}
\setlist{itemsep=3pt,topsep=5pt}
\setlength{\emergencystretch}{3em}
\pagestyle{fancy}
\fancyhf{}
\fancyhead[L]{\small\color{ink}Coupled computation and refinement in iPRA}
\fancyhead[R]{\small\thepage}
\renewcommand{\headrulewidth}{0.3pt}
\lstset{basicstyle=\ttfamily\small,breaklines=true,columns=fullflexible,
  keepspaces=true,backgroundcolor=\color{pale},frame=single,
  rulecolor=\color{lightgray},xleftmargin=5pt,xrightmargin=5pt}
\hypersetup{pdftitle={Coupled Computation and Refinement in iPRA},
  pdfauthor={Research report prepared for Vladimir Reshetnikov},
  pdfsubject={Constructive fixed points, residual certificates, and finite approximation towers}}

\begin{document}
\begin{titlepage}
\vspace*{14mm}
{\large\color{ink}MATHEMATICAL RESEARCH REPORT\par}
\vspace{11mm}
{\huge\bfseries\color{ink}Coupled Computation\\[5pt]and Refinement in iPRA\par}
\vspace{9mm}
{\Large Constructive fixed points, residual certificates,\\
and growing finite models\par}
\vspace{13mm}
{\large Prepared for Vladimir Reshetnikov\par}
\vspace{3mm}
19 November 2026\par
\vfill
\begin{minipage}{0.94\textwidth}
\textbf{Main result.} Finite rational contraction data produce a primitive
recursive regular sampler for a coupled system, even when the update is
available only through approximate implementations. An approximate
commutation estimate then makes finite models of different dimensions
cohere, with the contraction gap explicitly controlling the error.
\medskip

\textbf{Worked construction.} The report constructs and proves uniqueness
of a solution, in the specified box, to a countably infinite nonlinear
system using only bounded-dimensional dyadic computations at each query.
\medskip

\textbf{Verification status.} Complete human proofs are provided for the
main finite arithmetic and represented-object results. Eighteen supporting
arithmetic lemmas are compiled in Lean 4.32.0. The full
\texttt{Prf PRAAx} export is an explicitly identified integration outline,
not a completed machine-checked iPRA formalization.
\end{minipage}
\vspace{9mm}
{\small Repository baseline: ProveIt commit\par
\texttt{3ee850b4c2d650cf01263823fd5415963747c5d0}\par}
\end{titlepage}

\begin{abstract}
The existing ProveIt programme already contains extensive theories of
finite certificates, names, demand propagation, and proof composition.
This report extends the mathematical producers used by those theories.
A nonnegative rational matrix and positive weights certify contraction
of mutually dependent computations. A finite affine recurrence produces
the number of iterations and working precision needed for a regular
name, without first assuming an ideal real-valued map or a fixed point.
Finite Neumann sums produce componentwise residual enclosures. A second
theorem converts commutation defects between finite models into
coherence bounds; it also identifies why a shrinking contraction gap
cannot be ignored. Polynomial boxes provide a concrete certificate
frontend, and an infinite nonlinear lattice supplies a fully worked
example in which the dimension grows with the observation request.
Every infinite-object assertion is paired with its first-order
arithmetic reading, its primitive recursive producers, and the
quantifier-free invariants used in its proof. The report distinguishes
these results from arbitrary cyclic proof justification, singular
branch selection, and the unresolved effective bounds in the
Diophantine and prime-distribution packets.
\end{abstract}

\tableofcontents
\newpage
\section{Results and intended use}
\label{sec:overview}

The useful extension is a way to produce compatible mathematical
observations from a system of numerical dependencies. It applies to
simultaneous equations, approximate inverses, interacting analytic
constructions, and finite models whose size increases as more information
is requested. Its input is finite algebra and explicit error estimates;
its output is a primitive recursive sampler with a proved modulus.

There are two central statements. First, suppose a rational implementation
\(E_p\) maps a certified box into itself and satisfies, componentwise,
\[
 |E_p(x)-E_r(y)|\le A|x-y|+(2^{-p}+2^{-r})w,
 \qquad Aw\le(1-1/s)w,
 \tag{1}\label{eq:overview-action}
\]
where \(A\ge0\), \(w_i>0\), and \(s\ge1\) is an integer. The weighted
distance is \(d_w(x,y)=\max_i |x_i-y_i|/w_i\). A finite iteration with an
explicit precision schedule constructs \(Q(k)\) satisfying
\[
 d_w(Q(k),Q(l))\le2^{-k}+2^{-l}.
 \tag{2}\label{eq:overview-name}
\]
Moreover, its residual tends to zero at a supplied rate, and any other
small-residual computation agrees quantitatively with this one. There
is no assumption that a completed real fixed point has already been
constructed. The maps \(E_p\) may be discontinuous rounding procedures;
the mixed estimate in (1) is the property used in the proof.

Second, suppose finite-level updates \(F_n\) and restrictions \(P_n\)
satisfy a commutation estimate
\[
 d_n(F_nP_n z,P_nF_{n+1}z)\le\eta_n,
\]
and the actual returned points \(y_n\) have residual at most \(r_n\).
If \(F_n\) contracts by \(q_n<1\), then
\[
 d_n(y_n,P_ny_{n+1})\le
 \frac{r_n+\eta_n+r_{n+1}}{1-q_n}
 \tag{3}\label{eq:overview-refinement}
\]
for nonexpansive restrictions. An explicit summable bound on the right
produces a coherent family of observations. This supplies the missing
quantitative connection between solving a finite model and obtaining one
object across successive models.

\subsection{A concrete infinite-dimensional endpoint}

For \(i\in\N\), consider
\[
 x_i=\frac1{4(i+1)}+\frac{x_i^2}{8}+\frac{x_{i+1}}8,
 \qquad 0\le x_i\le1.                         \tag{4}\label{eq:overview-lattice}
\]
The report constructs a single primitive recursive rational function
\(X(i,k)\). To answer one query it uses only
\(i+k+3\) coordinates, \(i+k+7\) rounded iterations, and
\(i+k+11\) binary fractional digits. It proves
\[
 \begin{split}
 |X(i,k)-X(i,l)|&\le2^{-\min(k,l)-2},\\
 \left|X(i,k)-\frac1{4(i+1)}-
 \frac{X(i,k)^2}{8}-\frac{X(i+1,k)}8\right|&\le2^{-k-4}.
 \end{split}                                                  \tag{5}
\]
These are uniform arithmetic assertions, not a finite table of test
values. They express a unique solution to (4) within the specified box.
Uniqueness is not asserted among all bounded sequences.

\subsection{What is proved, and at what level}

\begin{longtable}{@{}>{\raggedright\arraybackslash}p{0.27\textwidth}>{\raggedright\arraybackslash}p{0.41\textwidth}>{\raggedright\arraybackslash}p{0.24\textwidth}@{}}
\toprule
Result & Mathematical endpoint & Evidence in this packet\\
\midrule
\endhead
Coupled approximation & Explicit residual bound and regular sampler from
(1) & Full human proof; arithmetic core compiled in Lean\\
Matrix certificates & Positive weights, componentwise supersolutions,
finite Neumann bounds & Full finite proofs and exact reference checks\\
Polynomial systems & Root names and parameter stability from a finite
box certificate & Full human proof and explicit coefficient algorithms\\
Refinement towers & Coherence from residuals, commutation defects and
contraction gaps & Full finite proof and formalization definitions\\
Infinite nonlinear system & PR coordinate sampler, equations and uniqueness
on \([0,1]^{\N}\) & Full human proof and exact reference computation\\
Object-theory export & \(\Gamma\vdash[\texttt{PRAAx}]\varphi\) for the
coded endpoints & Detailed invariant plan and an uncompiled adapter\\
\bottomrule
\end{longtable}

The human arguments identify the primitive recursive definitions and
quantifier-free induction formulas required for iPRA. The compiled
Lean arithmetic lemmas check important finite identities and bounds.
They do not, by themselves, certify the entire syntactic translation into
\texttt{Prf PRAAx}. This separation is part of the deliverable's exact
claim, rather than a condition left for the reader to discover.

The underlying contraction and verified-numerics ideas are established
mathematics; compare Rump's finite-system treatment and the radii-polynomial
method \cite{rump1994,radii}. The contribution here is the particular
finitary organization, the approximate-action and cross-level interfaces,
the explicit schedules, and their integration with the existing iPRA
programme. No claim of priority for the general contraction principle is
made.

\section{What this extension adds to the existing programme}
\label{sec:source-audit}

\subsection{Audit boundary and proof interface}

The source comparison is pinned to ProveIt commit
\nolinkurl{3ee850b4c2d650cf01263823fd5415963747c5d0}.
It uses the current unification plan and derivation inventory, the proof-packet
index, and the particular analytic and foundational chapters cited below.
It is a focused audit of the relevant interfaces, not a verification of every
research packet. The repository already contains many frameworks for bounded
certificates, finite observations, scoped witnesses, reverse precision
demands, and residual hypotheses. Merely giving those constructions another
name would not resolve a new mathematical difficulty.

The maintained interface is ordinary intuitionistic primitive recursive
arithmetic: intuitionistic first-order logic with equality, defining equations
for primitive recursive functions, successor axioms, and quantifier-free
induction. In the Lean development the public object judgment is
\(\Gamma\vdash[\texttt{PRAAx}]\varphi\).
The current unification plan expressly distinguishes an object derivation
from an ambient Lean theorem about standard natural numbers, and the
derivation inventory records the precise assumptions of each checked
endpoint \cite{audit-unification,audit-progress}.
Consequently a new rational estimate, a proof-producing arithmetic lemma,
and a completed analytic application have separate statuses throughout this
report.

\subsection{Scalar contractions and polynomial ODEs are established predecessors}

Two earlier analytic constructions should be credited explicitly.
The first embedding packet proves a finite scalar contraction estimate for
a normalized polynomial \(Q(X,Z)\), with contraction factor \(1/4\);
coefficient stabilization and a geometric tail then give a simple analytic
branch. Its perturbation estimate is independent of the number of iterations.
It also constructs a complex reciprocal from a finite derivative-separation
certificate, rather than assuming a reciprocal operation on arbitrary
Cauchy names \cite{audit-scalar}.

A subsequent chapter goes further. Given a root name \(b\), an inverse
derivative name \(w\), the identities
\[
 P(0,b)=0,\qquad wP_Y(0,b)=1,
\]
and numerical magnitude bounds, it chooses rational samples and small
radii by explicit formulas. It thereby \emph{produces} a passing
simple-root contraction certificate \cite{audit-scalar-producer}.
Thus the present work does not claim to introduce scalar contraction,
quantitative inversion, or the idea of producing an analytic chart.

Similarly, the local-to-global packet already gives a detailed human
construction of solutions to rational polynomial differential systems.
Dimension and degree are numerical parameters. Finite coefficient sums
produce bounds \(M,L\); a local horizon \(T=1/(2M)\) is constructed;
Euler polygons are compared through discrete Gronwall and adjacent-mesh
estimates; an explicit reindexing gives a regular solution name.
The same chapter proves uniqueness, perturbation estimates, and polynomial
first integrals \cite{audit-ode}.
Its perturbation discussion already accommodates rounded steps or
Cauchy-named coefficients when common domain and invariance contracts are
available.

The remaining formal boundary is narrower than the absence of an ODE
argument. The maintained progress inventory describes
\texttt{Foundations/DiscreteGronwall} as checked cleared arithmetic, and
\texttt{Foundations/OdeRefinement} as checked refinement arithmetic with the
error recurrence supplied. It keeps the complete rational polygon and
regular-name frontends separate \cite{audit-progress}.
A Picard or quadrature application of the present theorem is therefore an
alternative, more modular presentation of part of this analysis and an
extension to appropriately presented integrands. It is not the first
constructive polynomial ODE theorem in the repository.

\subsection{The specific gap: numerical feedback and a produced modulus}

The distinction between a cyclic precision calculation and execution of a
cyclic mathematical system is explicit in the current sources.
The checked two-coordinate demand example iterates
\[
 (x,y)\longmapsto
   \bigl(\min(B,y+1),\,\min(C,x+1)\bigr)
\]
and derives stabilization by the sum \(B+C\).
Its module documentation identifies it as a precision-demand instance, not
certificate execution \cite{audit-demand}.
The state-adaptive realizer implementation likewise handles bounded folds
using finite candidate tables and explicit paired selector alignment;
its inventory leaves broader cyclic execution and producer scopes
separate \cite{audit-adaptive,audit-progress}.

The new sufficient condition is quantitative and numerical.
For a finite family of state components, let \(A\) be a nonnegative rational
gain matrix. Supply positive rational weights \(w_i\) and a rational
\(q<1\) such that
\[
                     Aw\le q w
\]
coordinatewise. Local implementations have prescribed error budgets and
preserve a certified domain. The weighted maximum norm then turns the
coupled system into a contraction. A finite recurrence controls
iteration error and produces a Cauchy schedule, including the accumulation
of local approximation errors.

This addresses the feedback subclass for which the gain and domain
certificates hold. It does not justify circular proofs. Every local
inequality must already have an ordinary derivation; those derivations
are composed and used in a finite iteration argument. Nor does it solve
arbitrary recursive certificate systems. The numerical contraction is
additional mathematical information, and the theorem makes its use
visible.

This also strengthens the existing summable-motion interface in a useful
direction. The FIBER analysis chapter correctly distinguishes individually
valid approximations, finite stable windows, and a global Cauchy modulus.
Its telescoping theorem consumes a primitive recursive bound on every
finite tail \cite{audit-fiber}. Here the geometric tail is derived from
the gain certificate and the local error schedule. The contribution is the
\emph{producer of the tail and of a compatible sampler}; finite
telescoping and Booleanized induction are reused foundations.

\subsection{Approximate commutation extends the earlier tower methods}

There are already two distinct tower constructions in the repository.
The coherent-methods packet presents exact restrictions
\(\pi_{m,n}\) and proves coherence when a fixed producer satisfies a
specification preserved by restriction and unique at each finite level
\cite{audit-coherent}.
The topology framework uses finite simplicial complexes and refinement
carriers. Its map-coherence datum is naturality up to a common carrier,
hence contiguity, rather than equality of selected vertices
\cite{audit-topology}.
FINER's observational descent transports supplied agreement estimates
along a finite spanning tree \cite{audit-finer}.

None of these inspected chapters supplies the following analytic comparison
from a commutation defect. Let \(T_n,T_m\) be two finite-level updates,
\(R\) a restriction of Lipschitz constant \(L_R\), and suppose
\[
 \|RT_m z-T_n Rz\|\le\delta .
\]
Assume \(T_n\) contracts by \(q<1\), and the actual returned observations
satisfy
\[
 \|x_n-T_nx_n\|\le\varepsilon_n,\qquad
 \|x_m-T_mx_m\|\le\varepsilon_m .
\]
The triangle inequality and the contraction give the finite estimate
\[
 (1-q)\,\|x_n-Rx_m\|
       \le
       \varepsilon_n+L_R\varepsilon_m+\delta .
\]
Indeed insert \(T_nx_n\), \(T_nRx_m\), and \(RT_mx_m\)
between \(x_n\) and \(Rx_m\), bound the four differences, and move the
term \(q\|x_n-Rx_m\|\) to the left. This derivation uses only
finite rational inequalities; it does not assume exact fixed points.

An explicit bound on \(\delta\), together with selected numerical
precisions and a lower bound on \(1-q\), therefore converts finite-level
solvers into coherent finite observations. This is the precise new
connection with projective and spatial discretization.
The arithmetic target quantifies over finite level, observed coordinate,
iteration count, and demanded precision. An external interpretation as an
infinite sequence or inverse-limit object is justified by these contracts;
no function quantifier or compactness axiom is inserted into iPRA.

The denominator \(1-q\) is substantive. If the contraction factors
approach one, refinement defects must decrease sufficiently rapidly to
compensate, and the required schedule must be given explicitly.
Neither exact restriction coherence nor a useful limiting modulus follows
from small defects alone.

\subsection{Which further difficulties the theorem addresses}

The same finite comparison supports three reusable constructions.
A Neumann identity for finite matrix sums produces an inverse name and
a residual-to-solution error estimate whenever a rational
preconditioner has a certified contraction defect.
A multidimensional implicit system can use
\(T(x)=x-BF(x)\), with a finite box-invariance test and a bound on
\(I-BJ(x)\), to produce simultaneous root coordinates.
A finite quadrature implementation of a Picard map can use the same
error recurrence to control its discretization and evaluation errors.
Each application must derive its local bounds from its actual numerical
presentation. A claim that an inverse, derivative gap, or quadrature
modulus merely exists does not discharge that work.

These are extensions of the existing local scalar and ODE constructions
to coupled and approximately commuting computations.
They do not remove the need for separation where a problem asks for
a discontinuous label, or the need for a domain invariant on a longer
time interval. In particular, singular branch selection and global
existence beyond the certified region are not consequences of a local
contraction theorem.

\subsection{Diophantine and ineffective-constant boundaries remain distinct}

The Faltings packet distinguishes a positive height/list theorem from
a cardinality statement: a complete list and a height bound are equivalent
for its concrete point coding, but a bound on the number of points does not
tell an enumeration when to stop \cite{audit-faltings}.
This is not an assertion that every possible arithmetization of finiteness
must first produce a height bound.

The general-route chapter of the Siegel packet makes that distinction
especially clear. It proposes a finite-tuple covering theorem followed
by a zero-count theorem to rule out lists of too many distinct points.
That route could establish a cardinality conclusion without locating
all points. However, the finite approximation-cover theorem, the required
geometric presentations, and the cover-and-lift estimates are not supplied
by the finite counting combinator \cite{audit-siegel}.
The new contraction machinery can help evaluate local analytic data once
the needed quantitative contracts have been proved. It does not produce
these global arithmetic contracts.

The prime-distribution packet similarly isolates an effective
Siegel-type input needed for its requested uniform bounds, while proving
finite ingredients such as a weighted Vaughan identity and a rational Gram
certificate \cite{audit-primes}. Explicit local error budgets cannot turn
an ineffective existence proof into numerical constants.

Accordingly the resolution is partial but substantive: a broad class of
coupled and discretized constructions gains an actual primitive recursive
solver, a produced modulus, and a finite arithmetic proof route.
General Faltings, general Siegel, unrestricted branch choice, and the
missing effective prime-distribution inputs remain separate research
problems. Their unresolved status is not a proof of independence or
unprovability in iPRA.


\section{The arithmetic language and the meaning of an embedding}
\label{sec:arithmetic}

\subsection{The base theory}

The object theory in this report is first-order intuitionistic logic with
equality over natural numbers. Its nonlogical axioms are the defining
equations of primitive recursive functions, successor injectivity,
\(Sx\ne0\), and induction for quantifier-free formulas, with numerical
parameters. Logical rules are ordinary intuitionistic introduction and
elimination rules for \(\bot,\wedge,\vee,\to,\forall,\exists\), together
with equality reflexivity and substitution. There is no added choice
principle, unrestricted induction, excluded middle for arbitrary
formulas, or axiom of analytic completeness.

The presentation with a symbol for each primitive recursive definition is
a convenient definitional language. The repository's ordinary
\texttt{PR}, \texttt{PTm}, \texttt{Fm}, \texttt{PRAAx} and \texttt{Prf}
provide the intended formal syntax. Its bridges and scalar
bi-interpretation can be used after the ordinary derivation is obtained
\cite{audit-unification}. The report needs no blanket conservation claim
about a stronger analytic system. Foundational comparisons such as
\cite{fefermanstrahm} give useful context, but are not substituted for the
specific finite proofs below.

\subsection{Explicit finite coding}

One may fix the following coding throughout. Encode an integer \(z\) by
\(2z\) when \(z\ge0\), and by \(2(-z-1)+1\) when \(z<0\).
Encode a rational by a signed numerator and a natural \(d\), representing
the fraction with denominator \(d+1\). Reduced fractions are optional.
Rational equality and order are tested by signed cross multiplication;
addition and multiplication use the usual product denominators. All
operations are total primitive recursive functions on the codes.

For pairs use
\[
 \langle a,b\rangle=(a+b)^2+a.
\]
Different sums \(a+b\) occupy disjoint intervals, and the final summand
determines \(a\) within an interval. Both projections can therefore be
defined by bounded search through \(a,b\le c\), returning zero on unused
codes. The identities on actual pair codes follow from these finite
interval inequalities. A list is a length paired with a nested pair
chain; entry extraction performs the indicated number of projection
steps and returns zero outside its length. Thus dimension, matrix
entries, exponent vectors and iteration traces are numerical data.

A rational polynomial in \(d\) variables is a list of pairs consisting of
a coefficient and a length-\(d\) exponent vector. Its value is
\[
 \sum_{\alpha}c_{\alpha}\prod_{j<d}x_j^{\alpha_j}.
\]
Finite addition, multiplication, substitution and coefficient comparison
are primitive recursive on these data. This is a uniform evaluator for
\emph{this particular bounded syntax}, not a universal evaluator for
all primitive recursive programs.

The maximum of a nonempty finite list of rationals is computed by a
left-to-right comparison scan. At equal entries keep the first index.
This provides an actual attaining index whenever a proof below uses a
maximum. Absolute value, coordinate clipping to a rational interval,
and dyadic rounding are also decidable rational operations. Invalid
certificate inputs can be given a harmless default output; correctness
theorems are implications from the decidable validity predicate.

\subsection{Bounded formulas really become quantifier-free}

For a fixed primitive recursive Boolean predicate \(b(\vec u,j)\), define
\[
 \operatorname{All}_b(\vec u,0)=1,\qquad
 \operatorname{All}_b(\vec u,t+1)
   =\operatorname{All}_b(\vec u,t)b(\vec u,t).
\]
The two factors are bits. Quantifier-free induction proves that the
product is a bit and that it equals one exactly when every inspected
entry equals one. Bounded existential search is obtained by a first
successful-index scan with the bound as failure sentinel. Its invariant
records the current chosen index and whether a success has occurred.
Consequently any fixed bounded formula whose atomic predicates are
primitive recursive has a primitive recursive characteristic bit,
together with an intuitionistic proof of its equivalence to that bit
being one. The repository already supplies the corresponding
\texttt{BoundedFormulas} interface.

This is the precise justification for induction on statements such as
``all \(d\) current error coordinates obey their bounds.'' The actual
induction formula is one equality between numerical terms, optionally
combined with other quantifier-free equations and implications. It is
not an induction on a formula quantifying over all vectors or all
future precisions.

For example, if \(a(t)\) and \(b(t)\) are registered PR rational
sequences and \(C(t)\) is the bit checking the first \(t\) recurrence
inequalities, a usable induction formula is
\[
 C(t)=1\ \longrightarrow\ d^t a(t)\le b(t),
 \tag{6}\label{eq:qf-invariant}
\]
after denominator clearing. All variables other than \(t\) are numerical
parameters. At the successor step, \(C(t+1)=1\) supplies both
\(C(t)=1\) and the next recurrence inequality. This avoids silently
using induction for a universally quantified invariant.

\subsection{Registered producers and numerical parameters}

In theorems mentioning an evaluator \(E\), the object-language use is
schematic in one fixed PR definition of \(E\), with input parameters
\(u\), precision, and finite vector code as arguments. A polynomial
implementation can instead be uniform in its finite polynomial code,
because that restricted evaluator was just defined. Either version
allows dimension, coefficient numerators, denominator bounds, iteration
count and output precision to vary numerically.

The hypotheses about the evaluator are ordinary first-order formulas.
For example, ``\(E_p\) maps the box into itself'' means a universal
statement over \(p\) and vector codes with a bounded coordinate test.
It is used by specialization at the concrete current state. A generic
theorem retains that hypothesis. An application removes it by the
finite algebraic argument specific to its evaluator.

\subsection{Names, equality and equations}

\begin{definition}[Regular sampler]
A regular rational sampler in a weighted finite box is a fixed PR
function \(Q(k)\) such that all its outputs lie in the box and
\[
 d_w(Q(k),Q(l))\le2^{-k}+2^{-l}\quad(k,l\in\N).
\]
These inequalities, with denominators cleared, are the arithmetic
content. Two samplers describe the same value if their observations
agree to arbitrary requested precision after explicit PR reindexings.
\end{definition}

For a rational polynomial equation, a represented solution can be
specified by a regular sampler and an explicit residual estimate, for
example \(\abs{P(Q(k))}\le C2^{-k}\). Reindexing by a boundedly computed
binary magnitude of \(C\) converts this to any requested error. One
does not test equality of two arbitrary represented reals. The equality
is the family of rational estimates that has just been proved.

For a countable array the sampler has the form \(Q(i,k)\). The coordinate
index is a numerical input, not a second-order variable. For an externally
supplied name the proof is a schema for its fixed PR evaluator, or a
finite-query contract for tables containing exactly the requested
samples. Arbitrary external reals need not have such an evaluator.

\subsection{Small errors and exact rational conclusions}

A frequently used finite rule is worth making explicit. If rational
\(a>b\), write the positive difference as \(u/v\), \(u,v\in\N_{>0}\).
Then \(a-b\ge1/v\). Since \(2^v\ge v+1\), a bound
\(a-b\le2^{-v}\) is already contradictory. Thus, when a rational
inequality is proved with arbitrarily small dyadic slack, one particular
precision, computed from its denominators, removes the slack. This is
not a limit axiom. The same observation justifies the finite
cross-multiplication and error-removal steps used in represented
equations.

Finally, ``primitive recursive'' is a logical resource bound, not a
promise of numerical efficiency. Some exact rational iterates have
rapidly growing denominators. The rounded construction in
\cref{sec:lattice} is included precisely to keep the reference
computations modest while retaining exact error proofs.

\section{Coupled rational actions: producing convergence and error bounds}
\label{sec:small-gain}

The previously reviewed scalar implicit-branch constructions already turn
specific polynomial contractions into regular names.  The extension here is
a reusable \emph{matrix} interface: it handles mutually dependent coordinates,
precision-dependent rational actions, anisotropic error certificates, and
bounded-delay coordinate evaluation.  Its hypotheses and conclusions concern
only finite rational calculations.  No ideal map on a completed real space is
needed to construct the output.

\subsection{Finite input data and the approximate-action contract}

Fix a positive dimension $d$.  Let
$D=\prod_{i<d}[\ell_i,u_i]\subseteq\mathbb Q^d$ be a nonempty rational box,
with a specified seed $x_0\in D$.  Let $w_i>0$ be rational weights and put
\[
 |v|_w=\max_{i<d}\frac{|v_i|}{w_i},\qquad
 d_w(x,y)=|x-y|_w,\qquad e_p=2^{-p}.
\]
The maximum is the output of a finite scan.  For an integer $s\ge1$ put
$c=1-1/s$.  A small-gain certificate consists of a nonnegative rational
matrix $A$ and the finite inequalities
\begin{equation}
 Aw\le cw.                                      \tag{SG1}
\end{equation}
All vector inequalities are componentwise.  A rational number $R\ge0$
with $u_i-\ell_i\le Rw_i$ is a diameter certificate.

Let $E(x,p)$ be a \emph{fixed} primitive recursive rational-vector program.
Its input contracts are
\begin{align}
 x\in D&\Longrightarrow E(x,p)\in D,              \tag{SG2}\\
 |E(x,p)-E(y,q)|
 &\le A|x-y|+(e_p+e_q)w\quad(x,y\in D).           \tag{SG3}
\end{align}
These formulas are universally closed arithmetic contracts for the fixed
program, not a proposed universal evaluator for arbitrary primitive
recursive program codes.  In particular, putting $x=y$ in (SG3) supplies
the required coherence between work precisions.  From (SG1)--(SG3),
\begin{equation}
 d_w(E(x,p),E(y,q))\le c\,d_w(x,y)+e_p+e_q.       \tag{SG4}
\end{equation}
Indeed $|x-y|\le d_w(x,y)w$; multiply this inequality by $A\ge0$,
use (SG1), and divide each coordinate by its positive weight.

\subsection{A finite comparison principle}

\begin{lemma}[Weighted maximum principle]
\label{lem:weighted-maximum}
Suppose (SG1) holds.  For rational vectors $v,r,z$, if
\[
 v\le Av+r,\qquad z\ge Az+r,
\]
then $v\le z$.
\end{lemma}
\begin{proof}
Set $h=v-z$.  The two hypotheses give $h\le Ah$.
Compute $M=\max_{i<d}h_i/w_i$ and an index attaining the maximum.
If $M\le0$, the conclusion follows.  If $M>0$, choose an attaining
index $i$.  Since $h_j\le Mw_j$ and $A_{ij}\ge0$,
\[
 Mw_i=h_i\le(Ah)_i\le M(Aw)_i\le cMw_i.
\]
Thus $(1-c)Mw_i\le0$, contrary to $1-c=1/s>0$.
This is a decidable rational case split and a finite maximum, not a maximum
principle for functions on an infinite domain.
\end{proof}

Define the scalar and vector residuals by
\[
 r_p(x)=d_w(x,E(x,p))+e_p,\qquad
 \rho_p(x)=|x-E(x,p)|+e_pw.
\]
The triangle inequality and (SG3) give
\begin{equation}
 |x-y|\le A|x-y|+\rho_p(x)+\rho_q(y).             \tag{SG5}
\end{equation}
The scalar consequence of (SG5) is
\begin{equation}
 d_w(x,y)\le s\bigl(r_p(x)+r_q(y)\bigr).         \tag{SG6}
\end{equation}
For completeness, the triangle inequality first gives
$d_w(x,y)\le c\,d_w(x,y)+r_p(x)+r_q(y)$.
Multiply by $s$ and cancel $(s-1)d_w(x,y)$ on the two sides.
This produces an explicit isolation modulus from a finite matrix
certificate.  It does not ask the user to supply an unspecified uniqueness
modulus.

\subsection{A regular-name producer with an explicit schedule}

Choose natural numbers $H,G$ satisfying
\[
 sR\le2^H,\qquad s(2s+1)\le2^G.                 \tag{SG7}
\]
These can be found by bounded search: for a nonnegative rational $b$,
search $h\le\lceil b\rceil+1$ for $b\le2^h$.
The elementary inequality $h+1\le2^h$ proves the bound sufficient.
For fixed work precision $p$, define
\[
 x^p_0=x_0,\qquad x^p_{n+1}=E(x^p_n,p).
\]
For requested output precision $k$, set
\begin{equation}
 P(k)=k+G+1,\qquad N(k)=s(k+H+1),\qquad
 Q(k)=x^{P(k)}_{N(k)}.                           \tag{SG8}
\end{equation}

\begin{theorem}[Coupled approximate-action producer]
\label{thm:small-gain-producer}
Under (SG1)--(SG3) and the diameter certificate, the fixed primitive
recursive program (SG8) satisfies, for all $k,l$,
\begin{align}
 Q(k)&\in D,                                    \tag{SG9}\\
 s\,r_{P(k)}(Q(k))&\le e_k,                     \tag{SG10}\\
 d_w(Q(k),Q(l))&\le e_k+e_l.                    \tag{SG11}
\end{align}
It therefore produces a regular rational name.  Every other pair of fixed
programs $Z(k),P_Z(k)$ satisfying $Z(k)\in D$ and
$s r_{P_Z(k)}(Z(k))\le e_k$ names the same solution, with the explicit
cross-bound
\[
 d_w(Q(k),Z(l))\le e_k+e_l.
\]
Moreover, for every $k,p$,
\begin{equation}
 d_w(Q(k),E(Q(k),p))\le e_k/s+e_p.              \tag{SG12}
\end{equation}
Thus the fixed-point assertion itself has a finite requested-precision
meaning.
\end{theorem}
\begin{proof}
Invariance (SG9) follows by induction from (SG2).
Put $\delta_n=d_w(x^p_n,x^p_{n+1})$.
The diameter certificate gives $\delta_0\le R$, and (SG4) gives
$\delta_{n+1}\le c\delta_n+2e_p$.
Induction proves
\[
 \delta_n\le c^nR+2se_p:
\]
the step uses $cs=s-1$ to obtain
$c(c^nR+2se_p)+2e_p=c^{n+1}R+2se_p$.
Consequently
\begin{equation}
 r_p(x^p_n)\le c^nR+(2s+1)e_p.                 \tag{SG13}
\end{equation}

Here is a purely rational estimate for the contraction clock.
Induction on $n$ proves
$c^n(1+n/s)\le1$, because
\[
 c\bigl(1+(n+1)/s\bigr)
 =1+n/s-(n+1)/s^2\le1+n/s.
\]
Taking $n=s$ gives $c^s\le1/2$.  Induction on $a$ therefore gives
$c^{sa}\le2^{-a}$.  This proof includes $s=1$, when $c=0$.
Substitute $n=N(k)$ and $p=P(k)$ into (SG13).  By (SG7),
\[
 \begin{split}
 s r_{P(k)}(Q(k))
 &\le sR\,2^{-(k+H+1)}
       +s(2s+1)2^{-(k+G+1)}\\
 &\le e_k/2+e_k/2=e_k.
 \end{split}
\]
This proves (SG10).  Apply (SG6) to prove (SG11) and the uniqueness
cross-bound.  Finally, apply (SG3) with equal point inputs and precisions
$P(k),p$.  Its contribution $e_{P(k)}+e_p$, together with the direct
residual, is $r_{P(k)}(Q(k))+e_p\le e_k/s+e_p$.
\end{proof}

The construction proves regularity before invoking any interpretation by
real numbers.  If regular names are subsequently interpreted as real
vectors, (SG12) is exactly the fixed-point conclusion.  An arbitrary
primitive recursive sequence is not declared convergent, and an arbitrary
computable map is not declared an admissible input.

\subsection{Anisotropic a posteriori error and perturbation}

\begin{theorem}[A finite supersolution is an error certificate]
\label{thm:anisotropic-residual}
Let $x\in D$, let $p$ be a work precision, and let $z\ge0$ be a rational
vector satisfying the finite test
\begin{equation}
 z\ge Az+\rho_p(x).                             \tag{SG14}
\end{equation}
Then every output of (SG8) satisfies
\begin{equation}
 |x-Q(k)|\le z+e_kw.                            \tag{SG15}
\end{equation}
\end{theorem}
\begin{proof}
Put $v=|x-Q(k)|$.  Equation (SG5) gives
$v\le Av+\rho_p(x)+\rho_{P(k)}(Q(k))$.
Since (SG10) bounds the last vector by $(e_k/s)w$, the vector
$z+e_kw$ is a supersolution:
\[
 A(z+e_kw)+\rho_p(x)+(e_k/s)w
 \le z+e_k(c+1/s)w=z+e_kw.
\]
Apply Lemma~\ref{lem:weighted-maximum}.
\end{proof}

The same proof treats parameter variation.  Suppose $E,F$ share the box,
weights and small-gain matrix, and have a cross-contract
\[
 |E(x,p)-F(y,q)|\le A|x-y|+b+(e_p+e_q)w
\]
with a nonnegative rational vector $b$.  If $z\ge Az+b$, their produced
names satisfy
\begin{equation}
 |Q_E(k)-Q_F(l)|\le z+(e_k+e_l)w.               \tag{SG16}
\end{equation}
Indeed the vector of their differences satisfies the comparison inequality
with forcing $b+(e_k+e_l)w/s$; add $(e_k+e_l)w$ to the supersolution.
When $b=B|a-a'|$ for finite parameter data, this is an explicit modulus
for the solution's dependence on those parameters.

\subsection{Producing action contracts from polynomial data}

The approximate-action interface has a concrete finite source.  Work in the
centred box $|x_i|\le R_i$ and fix a finite list of multi-indices $\alpha$.
For every output coordinate, let $a_{i\alpha}(p)$ be fixed rational
coefficient programs with certified bounds
\[
 |a_{i\alpha}(p)|\le M_{i\alpha},\qquad
 |a_{i\alpha}(p)-a_{i\alpha}(q)|
       \le\beta_{i\alpha}(e_p+e_q).
\]
Set $F_i(x,p)=\sum_\alpha a_{i\alpha}(p)x^\alpha$, and calculate
\[
 V_i=\sum_\alpha\beta_{i\alpha}\prod_jR_j^{\alpha_j},\qquad
 A_{ij}=\sum_{\alpha:\alpha_j>0}
 M_{i\alpha}\alpha_jR_j^{\alpha_j-1}
                    \prod_{h\ne j}R_h^{\alpha_h}.
\]
An empty product is $1$, and the displayed restriction excludes negative
exponents even when a radius is zero.  Successively change the coordinates
in $x^\alpha$ to the coordinates of $y^\alpha$ and use
$u^n-v^n=(u-v)\sum_{h<n}u^hv^{n-1-h}$.  This proves the finite inequality
\[
 |F(x,p)-F(y,q)|\le A|x-y|+(e_p+e_q)V.
\]
If the rational test
\[
 \sum_\alpha M_{i\alpha}\prod_jR_j^{\alpha_j}\le R_i
 \quad(i<d)
\]
passes, every $F(\cdot,p)$ maps the box into itself.  Choose a natural $g$
with $V_i\le2^gw_i$ for every coordinate and define
$E(x,p)=F(x,p+g)$.  The preceding inequality becomes precisely (SG3),
and invariance gives (SG2).  The weights can be supplied directly or
constructed by (SG17).  These are polynomial coefficient, matrix, and
rational tests; no analytic mean-value theorem or real equality decision
is being imported.

A rational approximate inverse $B$ for a polynomial system $P(x)=0$ can
be used by expanding $F(x)=x-BP(x)$ about a rational centre.  Whenever
the finite box and small-gain tests pass, the theorem supplies a local
root and all the preceding error and perturbation estimates.  To conclude
$P=0$ from the fixed-point equation, include a finite inverse certificate
for $B$, or directly a left-inverse certificate $CB=I$.  Without that
certificate the fixed-point assertion only implies $BP=0$.
More explicitly, if the rational family is expanded exactly as
$E(x,p)=x-BP_{p+g}(x)$, then (SG12) and $CB=I$ give
\[
 |P_{p+g}(Q(k))|\le (e_k/s+e_p)|C|w,
\]
where $|C|$ is entrywise absolute value.  For requested coordinate error
$e_n$, compute $M=\max_i(|C|w)_i$, choose $a$ with $2M\le2^a$, and
take $k=p=n+a$.  The displayed residual is at most $e_n$.
For named coefficients, these finite inequalities, together with the
coefficient contracts, are the asserted equality $P=0$.  The construction
makes no decision about a singular Jacobian or the existence of a suitable
local box.

\subsection{Producing the certificates by finite matrix arithmetic}

One need not guess weights.  Suppose a supplied positive integer $m$ and
rational $0\le\theta<1$ pass the test
\begin{equation}
 A^m\mathbf1\le\theta\mathbf1.                 \tag{SG17}
\end{equation}
Set
\[
 w=\sum_{j<m}A^j\mathbf1,\qquad W=\max_iw_i,
 \qquad s=\max\left(1,\left\lceil\frac{W}{1-\theta}\right\rceil\right).
\]
Then $w_i\ge1$ and
\[
 Aw+\mathbf1=w+A^m\mathbf1,
 \qquad Aw+(1-\theta)\mathbf1\le w.
\]
Since $w_i/s\le W/s\le1-\theta$, this proves (SG1).
Only a finite power and finite sums occur.  The supplied value $m$, or a
separately proved bounded search for such a value, is essential; the rule
does not say to keep taking powers until success.

There is also a finite constructor for (SG14), without inverting a matrix.
For any nonnegative rational vector $r$ and any natural $n$, put
\begin{equation}
 u_n=\sum_{j<n}A^jr,\qquad
 t_n=|A^nr|_w,\qquad z_n=u_n+s t_nw.             \tag{SG18}
\end{equation}
The finite geometric identity is
$Au_n+r=u_n+A^nr$.  Therefore
\[
 Az_n+r
 =u_n+A^nr+s t_nAw
 \le u_n+t_nw+sct_nw=z_n.
\]
Taking $r=\rho_p(x)$ constructs an error certificate for every $n$.
The remainder has the explicit estimate
$s t_n\le s c^n|r|_w$, by induction using (SG1).
Increasing $n$ improves the certified truncation error at the known block
rate $c^s\le1/2$.  No matrix inverse, spectral-radius oracle, or infinite
Neumann series is a premise of this construction.

\subsection{A coupled example that fails the unweighted row test}

On $D=[0,4]\times[0,1]$ consider the rational map
\[
 T(x,y)=(1+2y,\ x/8+1/4),\qquad
 A=\begin{pmatrix}0&2\\1/8&0\end{pmatrix}.
\]
It maps $D$ into $D$, and its first row has sum $2$, so the unweighted
supremum norm does not certify contraction.  With $w=(4,1)$ and $s=2$,
however, $Aw=w/2$.  Its fixed point is $(2,1/2)$.
The approximate action can be obtained by rounding each coordinate of
$T(x,y)$ down to a multiple of $w_ie_p$; the result remains in $D$ and
has coordinate error at most $w_ie_p$.  Thus (SG3) holds exactly.
The diameter bound is $R=1$; the choices $H=1,G=4$ give the complete
schedule $N(k)=2(k+2)$ and $P(k)=k+5$.

For an exact-action residual demonstration, take
$x=(199/100,501/1000)$ and use
$r=|x-T(x)|=(3/250,9/4000)$.  The scalar weight certificate gives
$z_0=(3/125,3/500)=(0.024,0.006)$.
Formula (SG18) at $n=2$ gives
$z_2=(9/400,21/4000)=(0.0225,0.00525)$.
Both pass the exact rational supersolution test $z\ge Az+r$.
For comparison, the exact solution of $(I-A)z=r$ is
$(11/500,1/200)=(0.022,0.005)$, which also passes.  These are bounds
for the error of the supplied candidate, not assertions that the bounds
are attained.  The accompanying reference program checks every fraction.

\subsection{Bounded-delay coordinate evaluation}
\label{sec:async-small-gain}

The matrix interface also permits cyclic dependencies to be evaluated one
coordinate at a time.  A finite run consists of states $X^t\in D$.
At step $t$, a specified set of coordinates is updated.  For an updated
coordinate $i$, choose indices
\[
 \max(0,t-D_0)\le\tau(t,i,j)\le t\quad(j<d),
\]
form the mixed vector $v_j=X^{\tau(t,i,j)}_j$, and set
$X^{t+1}_i=E_i(v,p)$; other coordinates are retained.
The Cartesian-box hypothesis ensures that every mixed vector belongs to
$D$.  Let $B\ge1$ be an update-window bound, and put $L=B+D_0$.
It suffices that, in the first $B$ steps of each block
$[nL,(n+1)L)$ used by the run, every coordinate is updated at least once.
This is a decidable finite coverage certificate.

\begin{theorem}[A certified clock for delayed coordinate runs]
\label{thm:async-small-gain}
For any finite run with the stated delay and coverage certificates, any
$j$, and any block count $n$, one has
\begin{equation}
 d_w(X^{nL},Q(j))\le c^nR+se_p+e_j.            \tag{SG19}
\end{equation}
Consequently, choose $H_0,G_0$ with $R\le2^{H_0}$ and $s\le2^{G_0}$.
Run the certified schedule at precision
$p(k)=k+G_0+2$ for
\[
 T(k)=L\,s(k+H_0+2)
\]
steps, and call its output $Y(k)$.  Then $Y$ is a regular name for the same
solution as $Q$, with
\begin{equation}
 d_w(Y(k),Q(j))\le e_k/2+e_j.                  \tag{SG20}
\end{equation}
\end{theorem}
\begin{proof}
Fix $j$.  When a coordinate is updated from a mixed vector whose
coordinates are all within weighted error $M$ of $Q(j)$, (SG3) and the
residual estimate (SG10) bound its new weighted error by
\[
 cM+e_p+r_{P(j)}(Q(j))
 \le cM+e_p+e_j/s.
\]
Put $K=se_p+e_j$; then $cK+e_p+e_j/s=K$.
At the beginning of block $n$, use the invariant that every state in
the most recent delay window is within $c^nR+K$ of $Q(j)$.
It holds at block zero by the diameter bound, with any prehistory replaced
by the seed.  An update preserves the envelope and improves that coordinate
to $c^{n+1}R+K$.  After the first $B$ steps every coordinate has received
such an update.  During the next $D_0$ steps the weaker old-history
envelope still applies, so every further update also obeys the improved
bound.  At the end of the block all states in the last delay window are
past the first $B$ steps and hence obey the improved bound.  This proves
the block invariant and (SG19).

At the declared horizon, $c^{s(k+H_0+2)}R\le e_k/4$ and
$se_{p(k)}\le e_k/4$, proving (SG20).
To prove regularity without an infinitary limiting argument, for any
$k,l$ choose the finite index $j=k+l+2$.  The triangle inequality through
$Q(j)$ gives
\[
 d_w(Y(k),Y(l))\le(e_k+e_l)/2+2e_j\le e_k+e_l.
\]
Taking $j=k$ in (SG20) gives
$d_w(Y(k),Q(k))\le3e_k/2$, the usual regular-name equality bound.
\end{proof}

There is no assertion here about an arbitrary fair schedule.  The finite
coverage and delay bounds are the resources that produce the primitive
recursive clock.  They may be supplied by a periodic scheduler, or verified
for the concrete finite run.  Without them, a coordinate may be neglected
for an unbounded amount of time.

\subsection{Primitive recursion and the induction audit}

Encode signed rationals by an integer numerator and a positive natural
denominator; all comparisons are cross-products.  Finite vectors and
matrices are sequence codes with explicit dimensions.  Bounds such as
$Aw\le cw$ and membership in $D$ are Boolean folds over the dimension.
The maximum in Lemma~\ref{lem:weighted-maximum} is a left-to-right finite
scan returning both its value and a first attaining index.

Every displayed procedure is built from composition and primitive recursion:
matrix products and powers, finite sums, rational maxima, the bounded search
for $H,G$, the precision map, and the fixed-precision orbit.  The necessary
induction invariants are finite arithmetic formulas:
\begin{itemize}
\item for an orbit, $x_n\in D$ together with
      $d_w(x_n,x_{n+1})\le c^nR+2se_p$;
\item for a matrix power, $A^nw\le c^nw$ and the finite geometric identity;
\item for a Neumann certificate, the computed prefix identity
      $Au_n+r=u_n+A^nr$;
\item for an asynchronous block, the Boolean conjunction of the old-history
      envelope, the update mask, and the improved envelope for coordinates
      already updated in that block.
\end{itemize}
The requested precision, work precision, matrix data, comparison-name index,
and full finite horizon are held as numerical parameters during induction.
No invariant quantifies over all names, all precisions, or a completed
function space.  All bounded conjunctions are compiled to characteristic
primitive recursive predicates before quantifier-free induction is used.

To implement the results in the repository, first prove the rational and
sequence-code algebra in the ordinary \texttt{Prf PRAAx} interface.  Then
instantiate quantifier-free induction with the listed invariants and use
simultaneous cut to insert the concrete action contracts.  A host Lean proof
of the weighted inequalities is useful for checking the mathematics, but is
not itself the corresponding object-iPRA derivation.  The attached outline
records both boundaries explicitly.

\section{Two further producers: inverse matrices and multivariate charts}
\label{sec:applications}

The matrix machinery has uses beyond iterating a single update.
This section gives two explicit producers. The first constructs an
inverse name from a finite contraction certificate. The second
constructs a passing multivariate chart certificate from a regular
rational base point. Both use finite identities, with all precision
choices determined before execution.

\subsection{A finite Neumann inverse}

For positive weights \(w\), define the induced matrix bound
\[
 \norm{M}_w=\max_{i<d}\frac{\sum_{j<d}|M_{ij}|w_j}{w_i}.
\]
The triangle inequality is immediate from finite sums. Expanding a
matrix product gives
\begin{align*}
 \sum_j|(MN)_{ij}|w_j
 &\le\sum_t|M_{it}|\sum_j|N_{tj}|w_j\\
 &\le\norm{N}_w\sum_t|M_{it}|w_t,
\end{align*}
so \(\norm{MN}_w\le\norm M_w\norm N_w\). All sums are finite,
and reordering them is a finite distributive identity. Also
\(\norm I_w=1\).

Let \(K\) be a rational matrix with \(|K|\le A\), where
\(A,w,s\) pass the small-gain test and \(q=1-1/s\). Then
\(\norm K_w\le q\). Define
\[
 S_N=\sum_{j=0}^{N-1}K^j,\qquad S_0=0.
\]

\begin{theorem}[Constructed inverse near the identity]
\label{thm:inverse}
The finite matrices \(S_N\) satisfy
\begin{align}
 (I-K)S_N=S_N(I-K)&=I-K^N,                     \tag{M1}\\
 \norm{S_M-S_N}_w&\le s q^N\quad(M\ge N),    \tag{M2}\\
 \norm{S_N}_w&\le s.                          \tag{M3}
\end{align}
Choose \(H\) with \(s\le2^H\), and set
\[
 C(k)=S_{s(k+H+2)}.
\]
This is a PR regular matrix sampler. It represents a two-sided inverse
of \(I-K\), with an explicit residual rate.
\end{theorem}
\begin{proof}
Expand the products in (M1). Every intermediate power cancels in
the finite telescoping sum. Since \(K\) commutes with its own
powers, the same identity holds on both sides.
For (M2), submultiplicativity gives
\[
 \norm{S_M-S_N}_w\le\sum_{j=N}^{M-1}q^j
 =q^N\sum_{j=0}^{M-N-1}q^j\le s q^N.
\]
The last inequality follows without an infinite series: multiply
the finite sum by \(1-q=1/s\) to get \(1-q^{M-N}\le1\).
This also handles \(q=0\), with \(q^0=1\). Taking \(N=0\)
gives (M3). The block bound \(q^s\le1/2\) gives
\(s q^{s(k+H+2)}\le2^{-k-2}\), so (M2) proves regularity.
Finally (M1) bounds both inverse residuals by \(q^{s(k+H+2)}\).
The all-precision equations for a two-sided inverse therefore follow
from explicit finite inequalities.
\end{proof}

This matrix norm controls every entry:
\(\abs{M_{ij}}\le(w_i/w_j)\norm M_w\). To return entry
\((i,j)\) with the conventional dyadic error, shift the matrix
precision by a boundedly computed binary magnitude of \(w_i/w_j\).
Thus the matrix-valued notation entails an ordinary family of
rational coordinate names.

For a general rational matrix \(D\), a rational preconditioner
\(B\) can be used by setting \(K=I-BD\). If a rational
two-sided inverse \(B^{-1}\) is also supplied and checked, then
\(D=B^{-1}(I-K)\), and the constructed inverse is \(C B\).
At a finite cutoff, the two residuals are explicitly
\[
 (S_NB)D-I=-K^N,\qquad D(S_NB)-I=-B^{-1}K^NB.
\]
Their norm bounds are respectively \(q^N\) and
\(\norm{B^{-1}}_w\norm B_w q^N\). A binary magnitude of
\(\max(1,\norm{B^{-1}}_w\norm B_w)\) supplies the additional
precision shift for the right residual. Multiplication by \(B\)
in the regular sampler similarly uses a shift for \(\norm B_w\).
This formulation deliberately makes the preconditioner's invertibility
evidence visible. One may replace it with a finite determinant or
elimination proof, but should then include that proof rather than
silently cancel a possibly singular matrix.

\begin{proposition}[Stability of the inverse construction]
Suppose \(K,L\) both have weighted norm at most the same \(q<1\).
Their inverse names differ, in represented matrix norm, by at most
\(s^2\norm{K-L}_w\), where \(s=(1-q)^{-1}\) is an integer
as in the preceding construction.
\end{proposition}
\begin{proof}
For every positive integer \(j\), finite telescoping gives
\[
 K^j-L^j=\sum_{t=0}^{j-1}K^t(K-L)L^{j-1-t}.
\]
Summing through \(N-1\), the norm of the difference of the two
partial Neumann sums is at most
\[
 \norm{K-L}_w\sum_{a+b\le N-2}q^{a+b}
 \le\norm{K-L}_w
       \left(\sum_{a<N}q^a\right)
       \left(\sum_{b<N}q^b\right)
 \le s^2\norm{K-L}_w.
\]
The cases \(N=0,1\) are empty sums. Comparing different cutoffs
adds only the explicit tails (M2). At any requested precision,
choose the two cutoffs to make those tails smaller than the assigned
slack. This is the represented inequality asserted.
\end{proof}

\subsection{Automatic multivariate local charts at rational base points}

Let \(u\in\Q^d\) be unknown coordinates and \(v\in\Q^e\)
parameters, with \(d\ge1\). Let \(P(v,x)\in\Q[v,x]^d\) be
a finite polynomial system, and let \(c\in\Q^d\) satisfy the
coefficient-checkable equality \(P(0,c)=0\). Write
\(J=\partial_xP(0,c)\), computed by formal differentiation of
the polynomial table. Suppose rational matrices \(B,C\) pass
\[
 BJ=I,\qquad CB=I.                             \tag{M4}
\]
For example, when a nonzero rational determinant of \(J\) has
been checked, one can use its finite adjugate formula for \(B\)
and set \(C=J\). The identities are finite rational tests.

Expand the explicit polynomial
\[
 T(v,u)=u-BP(v,c+u)=
       \sum_{\alpha,\beta}t_{i,\alpha,\beta}u^\alpha v^\beta
       \quad\text{in coordinate }i.           \tag{M5}
\]
Equation (M4) and the base-point equality imply
\[
 t_{i,0,0}=0,\qquad t_{i,e_j,0}=0
 \quad(i,j<d).
\]
These are proved coefficient identities, not zero tests on
arbitrary represented values.

Compute nonnegative rational constants
\begin{align*}
 C_0&=\max_i\sum_{\alpha=0,\,|\beta|\ge1}
                       |t_{i,\alpha,\beta}|,\\
 C_1&=\max_i\sum_{|\alpha|=1,\,|\beta|\ge1}
                       |t_{i,\alpha,\beta}|,\\
 C_2&=\max_i\sum_{|\alpha|\ge2}
                       |\alpha|\,|t_{i,\alpha,\beta}|.
\end{align*}
All omitted sums are zero. Now set
\[
 r=\min\left(1,\frac1{4\max(1,C_2)}\right),
 \qquad
 h=\min\left(1,\frac{r}{2\max(1,C_0)},
                   \frac1{4\max(1,C_1)}\right).
 \label{eq:chart-radii}\tag{M6}
\]
Both radii are explicitly positive rationals.

\begin{theorem}[Produced multivariate chart certificate]
\label{thm:chart}
For every \(v\) with \(\norm v_\infty\le h\), the map
\(u\mapsto T(v,u)\) preserves \(\norm u_\infty\le r\) and
contracts by \(1/2\). Consequently it has a uniformly PR regular
solution sampler \(u(v,k)\) in that box, and
\(x(v,k)=c+u(v,k)\) represents a root of \(P(v,x)=0\).
This root is unique in \(\norm{x-c}_\infty\le r\).
The construction is uniform in the finite polynomial code, dimensions,
base point, and inverse certificate.
\end{theorem}
\begin{proof}
On the parameter box, the constant-in-\(u\) part of (M5) has
norm at most \(hC_0\le r/2\). For the difference between two
\(u\) inputs, factor each monomial difference coordinate by
coordinate. The sum of the absolute coefficients of the resulting
Lipschitz bound is at most
\[
 \sum_{\alpha,\beta}|\alpha|\,
 |t_{i,\alpha,\beta}|r^{|\alpha|-1}h^{|\beta|}.
\]
Terms with \(|\alpha|=0\) do not occur in this sum. The pure
linear terms vanish. The remaining linear terms are at most
\(hC_1\le1/4\); all terms with \(|\alpha|\ge2\) are at
most \(rC_2\le1/4\), because \(r,h\le1\).
More explicitly, take the componentwise gain entries
\[
 A_{ij}=\sum_{\alpha,\beta:\,|\alpha|\ge1}
 \alpha_j |t_{i,\alpha,\beta}|
 r^{|\alpha|-1}h^{|\beta|}.
\]
Monomial factorization gives the componentwise action inequality for
this matrix, and the row sums just proved are at most \(1/2\).
Thus \(A\one\le\one/2\), the literal small-gain certificate.
Comparing an input with zero proves
\(\norm{T(v,u)}_\infty\le r/2+\norm u_\infty/2\le r\).

The exact polynomial action, regarded as an approximate action
independent of its precision argument, satisfies the coupled
producer hypotheses with \(w=\one\), \(s=2\), and diameter
\(2r\). Apply \cref{thm:small-gain-producer}. Its finite
fixed-point residual is an error bound on \(BP(v,c+u)\).
Multiplication by the supplied \(C\) in (M4) gives an explicit
error bound on \(P(v,c+u)\) itself. For example, if
\(M=\max_i\sum_j|C_{ij}|\), a shift by a binary magnitude
of \(2M\) makes every coordinate residual at most the requested
dyadic error. All these operations are bounded rational operations.

For uniqueness, any represented root in the box is a represented
fixed point of (M5). Comparing its finite observations with the
constructed sampler uses the residual separation inequality (SG6).
The residuals can be made as small as desired by their explicit
moduli, yielding equality. No new existence assertion about
completed reals is needed.
\end{proof}

\subsection{Parameter dependence and named inputs}

Finite factorization in the parameter coordinates gives the computable
bound
\[
 K_v=\max_i\sum_{\alpha,\beta:,|\beta|\ge1}
   |\beta|\,|t_{i,\alpha,\beta}|
              r^{|\alpha|}h^{|\beta|-1}
\]
for the sup-norm Lipschitz constant in \(v\), on the certified
boxes. The perturbation theorem (SG16) therefore gives a solution
modulus with constant \(2K_v\). In elementary notation,
\[
 \norm{u(v)-u(v')}_\infty\le2K_v\norm{v-v'}_\infty.
\]
Its precise meaning is the corresponding inequality between
finite samples with their known dyadic slack.

If a parameter is supplied by a registered regular name, clip its
rational samples coordinatewise to the parameter box. Clipping
does not increase coordinate distance. Choose the working
parameter precision large enough that \(K_v\) times the input
error is at most the desired action error. Then the polynomial
coefficient construction in \cref{sec:small-gain} proves (SG3)
for the resulting rational evaluator. The same producer thus
handles those named parameters as well.

The theorem constructs local multivariate solution charts; a finite
collection of such charts can be compared using their explicit
uniqueness and parameter bounds. It does not infer that such charts
cover an arbitrary algebraic set, and it does not handle a singular
Jacobian without additional mathematical structure.

\section{Growing finite models and approximate commutation}
\label{sec:refinement}

Coupled contraction solves a fixed finite-dimensional problem. Many
mathematical constructions require another step: a solver on a small
model must agree with a solver on a larger model after restriction.
Exact commutation is often false because of boundary conditions,
truncation, or quadrature. This section derives the needed agreement
from a quantitative commutation defect.

\subsection{Finite-level data}

For each level \(n\ge1\), let \(D_n\) be an explicitly coded domain of
finite rational data, and let \(d_n\) be a rational-valued metric on it.
A pseudometric also suffices if equality of observations is understood
modulo zero distance. The domain predicate, metric, updates, restrictions, returned points
\(y_n\), and numerical bounds \(q_n,L_n,\eta_n,r_n\) are fixed PR families. Their algebraic properties are
proved locally or retained as explicit hypotheses.

Suppose
\[
 F_n:D_n\longrightarrow D_n,\qquad
 P_n:D_{n+1}\longrightarrow D_n
\]
satisfy, for rational \(0\le q_n<1\), \(L_n\ge0\), and
\(\eta_n\ge0\),
\begin{align}
 d_n(F_nx,F_ny)&\le q_n d_n(x,y),
                                    \label{eq:level-contraction}\\
 d_n(P_nu,P_nv)&\le L_n d_{n+1}(u,v),
                                    \label{eq:restriction-lipschitz}\\
 d_n(F_nP_nz,P_nF_{n+1}z)&\le\eta_n.
                                    \label{eq:commutation-defect}
\end{align}
Let the actual finite computations return \(y_n\in D_n\) with
\[
 d_n(y_n,F_ny_n)\le r_n.             \label{eq:level-residual}\tag{R1}
\]
In particular, \(y_n\) need not be an exact fixed point. It can be a
rounded iterate with a residual certificate checked after computation.

\begin{theorem}[Residual and commutation bound]
\label{thm:commutation}
For the preceding finite data,
\[
 d_n(y_n,P_ny_{n+1})\le e_n,
 \qquad
 e_n=\frac{r_n+\eta_n+L_nr_{n+1}}{1-q_n}.
 \label{eq:adjacent-error}\tag{R2}
\]
\end{theorem}
\begin{proof}
Set \(a=d_n(y_n,P_ny_{n+1})\). Insert the intermediate points
\(F_ny_n\), \(F_nP_ny_{n+1}\), and \(P_nF_{n+1}y_{n+1}\).
The four successive distances are bounded respectively by
\(r_n\), \(q_na\), \(\eta_n\), and \(L_nr_{n+1}\). Thus
\[
 a\le r_n+q_na+\eta_n+L_nr_{n+1}.
\]
Cancellation gives \((1-q_n)a\le r_n+\eta_n+L_nr_{n+1}\).
The denominator is a supplied positive rational, so division is an
ordinary rational operation. Equivalently, keep this last inequality
as the denominator-free endpoint.
\end{proof}

The result is asymmetric in the two contraction factors: it uses only
the lower-level factor \(q_n\). Reversing the comparison would require
a suitable lift and its own properties. No such lift is assumed here.

\subsection{Finite telescoping through restrictions}

Define \(P_{n,n}\) to be the identity and
\(P_{n,m+1}=P_{n,m}\circ P_m\). This fixes the orientation of every
restriction and gives coherence by definition. Put
\[
 L_{n,n}=1,\qquad L_{n,j}=\prod_{t=n}^{j-1}L_t\quad(j>n).
\]
For the iPRA proof, fix both endpoints and two input vector codes.
Run the two descending restriction trajectories in parallel, with the
current level \(m-t\). Induct on the rational inequality
\[
 d_{m-t}(z_t,z'_t)\le
 \left(\prod_{j=m-t}^{m-1}L_j\right)d_m(z_0,z'_0),
\]
with the two domain-membership bits. The step is one instance of (\ref{eq:restriction-lipschitz}).
All endpoints and input vectors remain numerical parameters. Thus the
finite composition has the asserted Lipschitz bound without an induction
on a universally quantified statement about all vectors.

\begin{theorem}[Finite refinement estimate]
\label{thm:finite-refinement}
For \(m\ge n\),
\[
 d_n(y_n,P_{n,m}y_m)
 \le\sum_{j=n}^{m-1}L_{n,j}e_j.
 \label{eq:refinement-sum}\tag{R3}
\]
In particular, if every \(L_j\le1\), then
\[
 d_n(y_n,P_{n,m}y_m)\le\sum_{j=n}^{m-1}e_j.
\]
\end{theorem}
\begin{proof}
The case \(m=n\) has zero distance and empty sum. For the successor
step, insert \(P_{n,m}y_m\), apply the induction hypothesis to the
first distance, and use the Lipschitz bound of \(P_{n,m}\) and (R2)
for the second distance. This adds precisely \(L_{n,m}e_m\).
The induction is on a finite rational inequality. Its coded form is
quantifier-free.
\end{proof}

This proof does not take an infinite sum. A tail modulus means a PR
function \(T(n)\) with a separately proved finite assertion
\[
 \sum_{j=n}^{m-1}e_j\le T(n)\quad(m\ge n).
 \label{eq:finite-tail}\tag{R4}
\]
For example, \(e_j\le C2^{-j}\) gives \(T(n)=2C2^{-n}\) by the
finite geometric identity. When restrictions expand, the required
bound applies instead to the weighted sum in (R3). A bound on the
unweighted errors would then be insufficient.

\subsection{Producing named observations}

An observation index \(i\) specifies a minimal level \(b(i)\).
For \(n\ge b(i)\), suppose a PR observer \(O_{n,i}:D_n\to\Q\)
satisfies
\begin{align*}
 |O_{n,i}(x)-O_{n,i}(y)|&\le B_i d_n(x,y),\\
 O_{n,i}(P_nz)&=O_{n+1,i}(z),
\end{align*}
with a known rational \(B_i\ge0\). Exact observational
compatibility can be replaced by an explicit adjacent error, as
described below.

\begin{theorem}[A coordinate sampler from a refinement tower]
\label{thm:tower-sampler}
Assume nonexpansive restrictions and (R4). Suppose a PR schedule
\(\nu(i,k)\ge b(i)\), nondecreasing in \(k\), satisfies
\[
 B_iT(\nu(i,k))\le2^{-k-2}.
\]
Then
\[
 X(i,k)=O_{\nu(i,k),i}(y_{\nu(i,k)})
\]
is a PR family of regular rational names. More precisely,
\[
 |X(i,k)-X(i,l)|\le2^{-\min(k,l)-2}.
 \label{eq:tower-name}\tag{R5}
\]
\end{theorem}
\begin{proof}
If \(k\le l\), set \(n=\nu(i,k)\), \(m=\nu(i,l)\).
The compatibility equations, composed finitely many times, identify
\(O_{m,i}(y_m)\) with \(O_{n,i}(P_{n,m}y_m)\). Now apply the
Lipschitz bound of the observer, (R3), (R4), and the schedule inequality.
The case \(l\le k\) is symmetric. Natural order is decidable, so
the case split is intuitionistically valid.
\end{proof}

If instead
\(\abs{O_{n,i}(P_nz)-O_{n+1,i}(z)}\le\zeta_{i,n}\), the adjacent
observation difference is at most \(B_ie_n+\zeta_{i,n}\). A finite
tail bound for that sum gives exactly the same argument. Thus observers
may themselves be approximations, provided their errors are budgeted.

\subsection{How residual budgets are chosen}

The factor \((1-q_n)^{-1}\) is a numerical resource. For example,
suppose \(L_n\le1\), positive rationals \(a_n\) have a known
summable tail, and the input data satisfy
\[
 \eta_n\le\frac{1-q_n}{2}a_n,
 \qquad
 r_n+r_{n+1}\le\frac{1-q_n}{2}a_n.
\]
Then \(e_n\le a_n\). The first condition is a property of the
models and their restrictions; the second can be enforced by running
the finite solvers accurately enough. It is not possible to repair an
arbitrary model defect merely by performing more iterations of each
solver. More iterations reduce \(r_n\), not \(\eta_n\).

When gaps vary, one conservative local choice for a desired error
majorant \(a_n\) is to choose each \(r_n\), for \(n>1\), below
both \((1-q_n)a_n/4\) and \((1-q_{n-1})a_{n-1}/4\), and to
require \(\eta_n\le(1-q_n)a_n/2\). Endpoints are handled using
only the available neighboring condition. These are finite positive
rational budgets, and the approximate-action producer gives a bounded
precision and iteration schedule for meeting them whenever its
local certificates are available.

\subsection{Why this is a useful embedding interface}

The theorem permits one to work entirely with finite models, then
derive coherent observations without assuming an inverse-limit element
or choosing one exact solution at every level. It is applicable to
truncated coupled recurrences, finite sections of linear systems, and
certified spatial or spectral approximations when their actual
commutation defects can be bounded. These are application patterns,
not a claim that every PDE or every spectral scheme automatically
supplies those estimates. The next section supplies all the data for
one infinite nonlinear example.

\section{A fully constructed infinite nonlinear system}
\label{sec:lattice}

We now eliminate every abstract domain, evaluator and tail premise in
one example. The equation is
\[
 x_i=a_i+\frac{x_i^2}{8}+\frac{x_{i+1}}8,
 \qquad a_i=\frac1{4(i+1)},\qquad 0\le x_i\le1
 \quad(i\in\N).
 \label{eq:lattice-system}\tag{L1}
\]
This is a genuinely infinite dependency: computing a coordinate by
recursion on its successors would not terminate. The construction
below instead proves a bounded amount of successor information
sufficient for any specified precision.

\subsection{Finite models and their certificates}

For \(n\ge1\), let \(D_n=([0,1]\cap\Q)^n\), with
\[
 d_n(x,y)=\max_{i<n}2^{-i}|x_i-y_i|.
\]
These are weighted maximum distances with weights \(w_i=2^i\).
Let \(P_n\) delete the last coordinate. Prefix restrictions are
nonexpansive and compose exactly.

Define \(F_n\) by
\[
 (F_nx)_i=a_i+\frac{x_i^2}{8}+\frac{x_{i+1}}8,
 \quad i<n,
 \qquad x_n:=0.
 \label{eq:lattice-finite}\tag{L2}
\]
All coordinates are nonnegative, and
\((F_nx)_i\le1/4+1/8+1/8=1/2\). Thus the map preserves the box,
with room for downward rounding.

For \(u,v\in[0,1]\),
\(\abs{u^2-v^2}=\abs{u-v}(u+v)\le2\abs{u-v}\).
Consequently
\[
 |(F_nx)_i-(F_ny)_i|
 \le\frac14|x_i-y_i|+\frac18|x_{i+1}-y_{i+1}|.
\]
The absent last-coordinate difference is zero. Multiplication by
\(2^{-i}\) bounds the two terms by \(d_n(x,y)/4\) and
\(d_n(x,y)/4\), respectively. Hence
\[
 d_n(F_nx,F_ny)\le\tfrac12d_n(x,y).             \tag{L3}
\]
The gain matrix has diagonal entries \(1/4\), superdiagonal entries
\(1/8\), and zeros elsewhere. With \(w_i=2^i\), every nonterminal
weighted row equals \(w_i/2\). This is a literal finite instance of
the matrix certificate.

The vectors \(F_nP_nz\) and \(P_nF_{n+1}z\) agree at every
coordinate except \(n-1\); there the difference is \(z_n/8\).
Thus
\[
 d_n(F_nP_nz,P_nF_{n+1}z)
 \le2^{-(n-1)}/8=2^{-n-2}.                     \tag{L4}
\]
This is the required commutation defect. It comes from the boundary
truncation, not from a numerical roundoff estimate.

\subsection{A bounded dyadic implementation}

For a nonnegative rational \(u\), define
\[
 \lfloor u\rfloor_p=2^{-p}\lfloor2^pu\rfloor.
\]
Integer division computes it, and
\(0\le u-\lfloor u\rfloor_p<2^{-p}\). Use
\emph{synchronous} iterations: all new coordinates are computed from
the previous vector. Set
\begin{align*}
 p_n&=n+8, & t_n&=n+4,\\
 v_{n,0}&=0, &
 v_{n,t+1,i}&=\left\lfloor(F_nv_{n,t})_i\right\rfloor_{p_n},\\
 y_n&=v_{n,t_n}.&&
\end{align*}
The rounded states lie in \([0,1/2]^n\). Each step has weighted
error at most \(h_n=2^{-p_n}\), since every factor \(2^{-i}\le1\).

Let \(z_{n,t}=F_n^t(0)\) be the exact rational iterates. These are
used for the proof; the reference implementation need not construct
their large denominators. With
\(\delta_t=d_n(v_{n,t},z_{n,t})\), (L3) gives
\[
 \delta_0=0,\qquad
 \delta_{t+1}\le\tfrac12\delta_t+h_n,
 \qquad \delta_t\le2h_n.                       \tag{L5}
\]
The last bound is a single quantifier-free induction.

The initial exact increment has distance
\(d_n(0,F_n0)=1/4\). Repeated use of (L3) bounds the exact residual
after \(t\) steps by \(2^{-t-2}\). Therefore
\begin{align*}
 d_n(y_n,F_ny_n)
 &\le d_n(y_n,z_{n,t_n})
       +d_n(z_{n,t_n},F_nz_{n,t_n})
       +d_n(F_nz_{n,t_n},F_ny_n)\\
 &\le3h_n+2^{-t_n-2}
   =7\,2^{-n-8}
   \le2^{-n-4}.
\end{align*}
This is a proved residual budget, not a stopping criterion based on
observed agreement of successive decimals.

\subsection{Coherence and the actual coordinate formula}

Apply \cref{thm:commutation} with \(q_n=1/2\), \(L_n=1\),
\(\eta_n=2^{-n-2}\), and \(r_n=2^{-n-4}\). It yields
\[
 d_n(y_n,P_ny_{n+1})
 \le2(2^{-n-4}+2^{-n-2}+2^{-n-5})
 =\frac{11}{16}2^{-n}<2^{-n}.                  \tag{L6}
\]
The finite geometric sum and \cref{thm:finite-refinement} imply
\[
 d_n(y_n,P_{n,m}y_m)\le2^{1-n}\quad(m\ge n).
 \tag{L7}
\]
Coordinate \(i\) is an observer at every level \(n>i\), with
Lipschitz bound \(2^i\). Define the output by the explicit formula
\[
 \boxed{X(i,k)=(y_{i+k+3})_i.}                 \label{eq:lattice-name}\tag{L8}
\]

\begin{theorem}[Uniform solution sampler]
\label{thm:lattice-name}
The function (L8) is primitive recursive and takes values in \([0,1]\).
For every \(i,k,l\in\N\),
\begin{align}
 |X(i,k)-X(i,l)|&\le2^{-\min(k,l)-2},          \tag{L9}\\
 \left|X(i,k)-a_i-\frac{X(i,k)^2}{8}
                      -\frac{X(i+1,k)}8\right|
 &\le2^{-k-4}.                               \tag{L10}
\end{align}
Thus it represents a solution of (L1) by entirely arithmetic contracts.
\end{theorem}
\begin{proof}
Every operation in (L8) is a bounded list operation, integer division,
or a PR iteration with the explicit bounds \(n=i+k+3\),
\(t=n+4\), \(p=n+8\). Membership in the box has already been
proved. If \(k\le l\), (L7) gives
\[
 |X(i,k)-X(i,l)|\le2^i2^{1-(i+k+3)}=2^{-k-2}.
\]
The reverse order gives (L9).

For (L10), put \(n=i+k+3\). In particular \(i+1<n\), so
the \(i\)-th equation of \(F_n\) uses the actual next coordinate,
not the zero boundary. Its residual contributes at most
\(2^ir_n\). The value \(X(i+1,k)\), however, is taken at level
\(n+1\). Replacing \((y_n)_{i+1}\) by \((y_{n+1})_{i+1}\)
contributes at most \(2^{i+1}2^{-n}/8\), by (L6). Therefore
\[
 \left|X(i,k)-a_i-\frac{X(i,k)^2}{8}-\frac{X(i+1,k)}8\right|
 \le2^{i-n-4}+2^{i-n-2}
 =5\,2^{-k-7}<2^{-k-4}.
\]
This also explains why the difference between successive model
dimensions must be included in the equation proof.
\end{proof}

\subsection{Uniqueness in the specified box}

Call a candidate admissible if it is a regular rational array \(U(i,p)\)
with \(0\le U(i,p)\le1\), and its same-precision residual satisfies
\[
 \left|U(i,p)-a_i-\frac{U(i,p)^2}{8}
                     -\frac{U(i+1,p)}8\right|\le2^{-p}
 \quad(i,p\in\N).                              \tag{L11}
\]
This is one useful normalized representation. Ordinary regular
solution names with a known residual modulus can be reindexed to
this format. No algorithm for extracting such a modulus from an
unspecified external function is assumed.

\begin{theorem}[Finite uniqueness certificate]
\label{thm:lattice-unique}
If \(U,V\) are admissible, then, for every \(i,k\), putting
\(p=i+k+5\) gives
\[
 |U(i,p)-V(i,p)|\le2^{-k-2}.                    \tag{L12}
\]
Consequently their represented coordinates are equal. In particular,
(L8) is the unique represented solution in this class.
\end{theorem}
\begin{proof}
Let \(N=i+k+3\), and form the first \(N\) coordinates of
\(U(-,p)\) and \(V(-,p)\), denoted \(u,v\). For each of the
first \(N-1\) equations their finite-map residual is at most
\(2^{-p}\). The last equation has the additional omitted boundary
coordinate, contributing at most \(2^{-N-2}\) in weighted distance.
Thus both finite-map residuals are at most
\(2^{-p}+2^{-N-2}\). The triangle inequality and (L3) give
\[
 d_N(u,v)\le2(2^{-p}+2^{-N-2})+\tfrac12d_N(u,v),
\]
and hence \(d_N(u,v)\le4\,2^{-p}+2^{-N}\). Multiplication by
\(2^i\) yields
\[
 |U(i,p)-V(i,p)|
 \le2^{i+2-p}+2^{i-N}
 =2^{-k-3}+2^{-k-3}=2^{-k-2}.
\]
Only coordinates through \(N\) and one finite precision are queried.
For an entirely explicit equality bound at a fixed shift, apply
(L12) with target index \(k+2\), so the common work precision is
\(p=i+k+7\). Regularity of each candidate, between \(k+2\)
and \(p\), then gives
\[
 |U(i,k+2)-V(i,k+2)|
 \le 2\bigl(2^{-k-2}+2^{-i-k-7}\bigr)+2^{-k-4}
 \le\frac{37}{64}2^{-k}<2^{-k}.
\]
Thus equality itself has a supplied finite precision schedule. The
constructed sampler satisfies the stronger residual (L10), so it is
admissible.
\end{proof}

\subsection{The box is mathematically necessary}

Uniqueness among \emph{all bounded} solutions would be false. The
following independent finite construction produces another solution,
with every coordinate in \([6,8]\).

\paragraph{A completely finite upper-branch construction.}
For $z\in[6,8]\cap\Q$, put
$R_i(z)=16-z-2/(i+1)\in[6,10]$.
Starting with $[2,4]$, perform $p+4$ rational bisections, retaining an
interval $[l,u]$ with $l^2\le R_i(z)\le u^2$; at each step test the
square of the midpoint.  Define $b_i(z,p)$ to be the final midpoint and
$E_i(z,p)=4+b_i(z,p)$.  The final interval width is $2^{-p-3}$.
Since $2\le l\le b_i\le u\le4$,
\[
 |b_i(z,p)^2-R_i(z)|\le 8\,(u-l)/2=2^{-p-1}.
\]
The exact identities for differences of squares give, entirely over
rationals,
\[
 \begin{split}
 6\le E_i(z,p)&\le8,\\
 |E_i(z,p)-E_i(z',q)|
 &\le\tfrac14|z-z'|+\tfrac18(2^{-p}+2^{-q}),\\
 \left|E_i(z,p)-a_i-E_i(z,p)^2/8-z/8\right|
 &\le2^{-p}/16.
 \end{split}
\]
For the second inequality, factor the difference of the two midpoint
squares and use that their sum is at least four.  For the third, multiply
the expression by eight; it becomes $R_i(z)-b_i(z,p)^2$.

Define the finite backward chain by
$x_N^{N,p}=7$ and $x_j^{N,p}=E_j(x_{j+1}^{N,p},p)$ for $j<N$.
For $M\ge N$ and $i<N$, its cross-precision estimate is
\[
 |x_i^{N,p}-x_i^{M,q}|
 \le4^{-(N-i)}+\tfrac16(2^{-p}+2^{-q}).
\]
Indeed the terminal discrepancy at coordinate $N$ is at most one, and
finite iteration of the preceding contraction inequality gives the first
term and a geometric sum bounded by
$(2^{-p}+2^{-q})(1/8)(4/3)$.

Now define the actual rational sampler
\[
 Y(i,k)=x_i^{\,i+k+3,\,k+4}.
\]
It uses exactly $k+3$ backward applications, each with $k+8$ bisections.
For $k\le l$, the cross-precision estimate gives
\[
 |Y(i,k)-Y(i,l)|
 \le4^{-k-3}+\frac{2^{-k-4}+2^{-l-4}}6
 \le\frac7{192}2^{-k}<2^{-k-2}.
\]
Thus $Y$ is a regular PR name with every sample in $[6,8]$.
To verify the original equation, the internal next coordinate of
$x^{i+k+3,k+4}$ differs from $Y(i+1,k)$ by at most
$4^{-k-2}+2^{-k-4}/3$.  Adding its contribution divided by eight
to the local residual gives
\[
 \left|Y(i,k)-a_i-Y(i,k)^2/8-Y(i+1,k)/8\right|
 \le\frac{2^{-k-4}}{16}
    +\frac{4^{-k-2}+2^{-k-4}/3}{8}
 \le\frac{11}{768}2^{-k}<2^{-k-6}.
\]
This constructs a separate bounded upper solution with explicit finite
contracts.  No completed real function or unspecified square-root routine
was used to define it.  Both the bisection invariant and the finite backward
error invariant are rational inequalities suitable for quantifier-free
induction.

\subsection{The exact first-order endpoint}

Let \(\mathsf{lat}(i,k)\) be the registered PR rational code defined
by (L8). Then the main iPRA targets are the universal closures of:
\[
 \begin{gathered}
 \mathsf{inUnit}(\mathsf{lat}(i,k))=1,\\
 \mathsf{ratLe}\bigl(
   |\mathsf{lat}(i,k)-\mathsf{lat}(i,l)|,
   2^{-\min(k,l)-2}\bigr)=1,\\
 \mathsf{ratLe}\bigl(
   |\mathsf{lat}(i,k)-a_i-\mathsf{lat}(i,k)^2/8
                            -\mathsf{lat}(i+1,k)/8|,
   2^{-k-4}\bigr)=1.
 \end{gathered}
\]
Here vertical bars and fractions abbreviate the explicit rational
code operations from \cref{sec:arithmetic}. Each displayed matrix
is quantifier-free after this notation is expanded. The universal
closures are logical introductions after the finite invariant
proofs, not additional induction axioms.

\section{Exact boundaries and counterexamples}
\label{sec:boundaries}

The hypotheses identify useful mathematical structure. Several of them
cannot be discarded merely because each finite computation terminates.
The following elementary examples make the boundaries testable.

\subsection{Separately contracting implementations can disagree}

On \([0,1]\cap\Q\), let \(E_p(x)=0\) for even \(p\) and
\(E_p(x)=1\) for odd \(p\). Each map is a constant, hence a
contraction of factor zero with zero fixed-point residual. Their
fixed points alternate. No regular name follows.

The mixed-precision estimate (SG3) fails: for large adjacent
precisions its right side at equal input is \(2^{-p}+2^{-p-1}\),
while its left side is one. Thus cross-precision coherence is a
separate input property, and cannot be inferred from contraction of
each implementation in isolation.

\subsection{Vanishing commutation defects do not overcome a closing gap}

Take \(D_n=[0,1]\cap\Q\), identity restrictions,
\(\eps_n=2^{-n-2}\), and
\[
 F_n(x)=(1-\eps_n)x+\eps_nc_n,
 \qquad
 c_n=\begin{cases}0&n\text{ even},\\1&n\text{ odd}.
 \end{cases}
\]
Each map has the exact rational fixed point \(y_n=c_n\), and
contraction factor \(q_n=1-\eps_n<1\). Direct subtraction gives
\[
 \sup_{0\le x\le1}|F_n(x)-F_{n+1}(x)|\le\eps_n.
\]
For example, when \(c_n=0,c_{n+1}=1\), the difference is
\(-\eps_nx/2-\eps_n/2\); the other parity is analogous.
The commutation defect tends to zero with an explicit dyadic rate,
but the exact fixed points alternate between zero and one.

There is no contradiction with (R2): the defect divided by the
contraction gap is \(\eps_n/(1-q_n)=1\). It does not have a
vanishing tail. This example shows why any theorem that omits the gap
factor would be false even for one-dimensional affine maps.

\subsection{A local numerical certificate is not a circular proof rule}

A system of equations may have cycles in its dependency graph. The
present method handles a subclass by replacing an implicit solution
with a finite sequence of explicit updates and proving an error
invariant for that sequence. The proof of each update estimate is
already available before the iteration theorem is applied.

This does not license a proof graph in which a claim assumes itself
through another unproved claim. No equation such as ``certificate
\(c_1\) is valid because \(c_2\) is valid, and conversely'' is
accepted merely because its dataflow diagram is finite. An unrelated
cyclic certificate system would need a sound induction, descent or
guardedness argument of its own. Thus the open cyclic-execution item
in the repository is narrowed for numerical feedback with small-gain
data, not solved in its full generality.

\subsection{No total primitive recursive universal evaluator}

Suppose a total PR function \(U(e,n)\) evaluated every unary PR
program at its code \(e\). Then \(g(n)=U(n,n)+1\) would itself be
PR, with some code \(a\). The defining property at \((a,a)\)
would give
\(U(a,a)=g(a)=U(a,a)+1\), a contradiction. Consequently
``all producer codes are PR'' cannot be implemented by adding a
single total PR evaluator for them.

Our two legitimate uses of uniformity are different. A fixed PR
body can have arbitrarily many numerical parameters. Alternatively,
a fixed bounded syntax such as finite rational polynomials has its
own PR evaluator. General program traces can also be checked with a
supplied fuel bound, but the claim that the bound is sufficient is
then an explicit obligation.

\subsection{Metastability and convergence remain different outputs}

A finite stable window does not by itself supply a tail modulus or
a coherent name. The proof-mining literature explicitly distinguishes
rates of convergence, finite learning information, and rates of
metastability \cite{kohlenbachsafarik}. The current construction obtains
a full modulus because the gap and the error schedule give a finite
geometric estimate. A theorem lacking those data may still have
valuable metastable consequences; it must retain that weaker output
instead of applying the regular-name theorem without its hypotheses.

\subsection{Cardinality is not a complete enumeration}

Even a bound of one possible witness does not generally produce a
search cutoff. For a program code \(e\), the relation ``\(t\) is
the first halting time of program \(e\)'' is decidable by bounded
simulation and has at most one witness. A uniform computable bound
containing that witness whenever it exists would decide halting by
simulation up to the returned bound. The usual diagonal argument
rules this out.

This elementary example explains the logical distinction; it is not
a reduction proving anything about the particular curves in the
Faltings or Siegel packets. Their requested height bounds, complete
lists, cardinality statements and finite-tuple obstructions are
different arithmetical targets. A new proof of one target would
need its own exact hypotheses and argument. No inference of
independence, unprovability in iPRA, or impossibility of every
constructive approach is made here.

\subsection{Failure of a certificate is inconclusive}

A failed weighted-row test does not prove nonexistence of a solution.
Weights, a smaller box, a different preconditioner, or another
mathematical method may succeed. Likewise, a passing local test does
not establish global continuation, uniqueness outside its region,
or behavior at a singular parameter. The numerical inequalities
state the scope of the result exactly.

\section{The checked arithmetic kernel and its iPRA translation}
\label{sec:checked-kernel}

The accompanying file \texttt{AffineBounds.lean} was compiled with Lean
4.32.0 using only \texttt{Std}.  All eighteen theorem declarations were
checked; their axiom reports contain only \texttt{propext} and
\texttt{Quot.sound}, or the smaller subset \texttt{propext}.  They contain
no \texttt{sorryAx} and do not depend on \texttt{Classical.choice}.
These are host-Lean arithmetic theorems.  Their proof status must not be
confused with an object-level theorem of the form
\(\Gamma\vdash_{\mathrm{iPRA}} A\).  The latter is the endpoint of the
explicit adapter described below, whose arithmetic obligations remain to
be filled with object-level derivations.

\subsection{Exact finite bounds}

For natural numbers \(c,d,b\) and a fixed forcing sequence \(r_n\), define
by primitive recursion
\[
 B_0=b,\qquad B_{n+1}=cB_n+d^n r_n.
\]
If \(d a_{n+1}\le c a_n+r_n\) and \(a_0=b\), then
\[
 d^n a_n\le B_n.                                      \tag{K1}
\]
Indeed, multiplying the step inequality by \(d^n\) and applying the
induction hypothesis gives
\[
 d^{n+1}a_{n+1}\le c(d^na_n)+d^n r_n
                 \le cB_n+d^nr_n=B_{n+1}.
\]
This is a quantifier-free induction on the displayed inequality.
For constant forcing \(r\), put
\[
 G_0=0,\qquad G_{n+1}=cG_n+d^n.
\]
One more quantifier-free induction proves
\(B_n=c^nb+rG_n\).  If \(d=c+1\), then
\[
 G_n+c^n=d^n.                                         \tag{K2}
\]
At the successor step the left side is
\(cG_n+d^n+c^{n+1}=c(G_n+c^n)+d^n=d^{n+1}\).
Consequently (K1) yields the subtraction-free estimate
\[
 d^na_n+rc^n\le c^na_0+rd^n.                          \tag{K3}
\]
All expressions in (K1)--(K3) are natural-number terms.  A rational
application must first clear denominators at the finite query in question;
these integer inequalities cannot be applied to unencoded real errors.

The clock is also explicit.  Induction gives
\[
 c^{n+1}+(n+1)c^n\le(c+1)^{n+1}.
\]
Taking \(n=c\) gives \(2c^{c+1}\le(c+1)^{c+1}\), including
\(c=0\). Multiplication of this block inequality proves
\[
 2^m c^{(c+1)m}\le(c+1)^{(c+1)m}.                     \tag{K4}
\]
Thus \(m\) blocks of \(c+1\) steps supply \(m\) binary digits of
geometric decay.  Neither logarithms nor a search for a stopping time is
needed.

\subsection{Simultaneous components}

Let \(M=(M_{ij})\) be a nonnegative integer matrix, let \(w_i\) be
nonnegative integer weights, and suppose the finite row certificates are
\[
 \sum_{j<N}M_{ij}w_j\le c w_i\qquad(i<N).
\]
Positivity of the weights is needed when interpreting a weighted norm,
but the following cleared inequalities themselves allow zero weights.
If
\[
 e_{0,i}\le w_i b,\qquad
 d e_{n+1,i}\le\sum_{j<N}M_{ij}e_{n,j}+w_i r_n,
\]
then the checked theorem \texttt{coupled\_envelope} states
\[
 d^n e_{n,i}\le w_iB_n\qquad(i<N).                    \tag{K5}
\]
For the induction, multiply the recurrence by \(d^n\), distribute this
factor through the finite sum, use the component induction hypotheses,
factor out \(B_n\), and apply the row certificate.  The result is
\(d^{n+1}e_{n+1,i}\le w_i(cB_n+d^nr_n)\).
The single invariant for the object calculus is the bounded conjunction
of (K5) for \(i<N\); the bounded-formula compiler turns it into one bit
equation before induction.

The module also checks a useful posterior certificate.  From
\(dv\le cv+r\) and \(c\le d\), one obtains
\((d-c)v\le r\), because
\(dv=(d-c)v+cv\) and additive cancellation is constructive.
For two components, if
\[
 Ax\le by+r,\qquad By\le cx+s,\qquad bc\le AB,
\]
then
\[
 (AB-bc)x\le Br+bs.                                  \tag{K6}
\]
Multiply the first inequality by \(B\), substitute the second after
multiplying it by \(b\), and apply the preceding cancellation lemma.
A strictly positive determinant gap makes (K6) a numerical residual
bound.  Its positivity is explicit certificate data, not an implicit
matrix-invertibility assumption.

\subsection{Exact object-calculus interface}

At ProveIt commit
\texttt{3ee850b4c2d650cf01263823fd5415963747c5d0}, the namespace
\texttt{LeanProofs.IPRA.Foundations.GuardedLoop} already provides a
structure \texttt{Loop}. Its fields are an actual proof of
quantifier-freeness and actual derivations of base, bounded transition,
and final implication in \texttt{Prf PRAAx}. The theorem
\texttt{loop\_at} exports the result at the supplied primitive-recursive
fuel. The namespace \texttt{TraceReflection} supplies \texttt{trace} for
the universally closed result. The namespace \texttt{BoundedFormulas}
supplies the syntax \texttt{BF}, the formula interpretation
\texttt{fmAt}, the primitive-recursive bit interpretation \texttt{bitAt},
and the object-level equivalence \texttt{bf\_reflect}.

The accompanying \texttt{ObjectAdapter.lean} records this interface using
the actual inspected names and types. It is an uncompiled integration
outline. To complete the instance, encode the state and its rational
denominators, define the bounded invariant, prove its base and transition
using finite-sum algebra and (K1)--(K6), and supply those proofs to
\texttt{Loop}. A proof in the Lean proposition \(a\le b\) cannot fill a
field expecting \(\Gamma\vdash[\mathrm{PRAAx}]\mathrm{le}(a,b)\).
The host proof is a guide to these derivations, not an unproved
reflection rule.

Finally, every function argument in the host-Lean statements is
schematic.  An object-level instance fixes primitive-recursive program
bodies or supplies explicit bounded evaluation certificates.  The
framework does not postulate a total primitive-recursive universal
evaluator, nor does the finite error estimate alone assert existence of
an infinite real fixed point; represented-object adequacy and the
necessary modulus are separate obligations.

\clearpage
\section{Lean definitions and a path to complete object formalization}
\label{sec:formalization}

The formalization has three different products. The distinction prevents
the common error of replacing an object-theory proof by a theorem of
the surrounding prover.

\begin{enumerate}
\item \texttt{AffineBounds.lean} contains eighteen proved theorems in
Lean about natural-number recurrences and finite sums. It is compiled.
\item The blueprint modules contain compiled definitions of rational
data, algorithms, guards and named target propositions. Defining a
proposition does not prove it. These modules introduce no new axiom
asserting their targets.
\item \texttt{ObjectAdapter.lean} is an uncompiled integration outline
for actual \texttt{Prf PRAAx} derivations, at the pinned ProveIt
revision. Its proof fields must be populated by ordinary iPRA proof
objects.
\end{enumerate}

The arithmetic proofs can be checked independently of Mathlib and
ProveIt. The final object export requires the existing repository
arithmetic library and the additional instances described here.

\subsection{Finite rational definitions}

The coupled blueprint uses the host definitions
\begin{lstlisting}
abbrev Vec (d : Nat) := Fin d -> Rat
abbrev Mat (d : Nat) := Fin d -> Fin d -> Rat

structure Data (d : Nat) where
  matrix : Mat d
  weights : Vec d
  reciprocalGap : Nat

def contraction (c : Data d) : Rat :=
  1 - 1 / (c.reciprocalGap : Rat)
\end{lstlisting}

It defines finite sums and maxima by folds, and defines
\texttt{weightedNorm}, \texttt{InBox}, \texttt{SmallGain},
\texttt{Diameter} and \texttt{ActionContract} exactly as in the
human theorem. The precision and iteration functions are literal
natural-number expressions, and \texttt{orbit} is a primitive
recursion. Other definitions include vector residuals, finite Neumann
prefixes, automatic weights, finite state histories, delay guards,
and update-coverage guards.

The named propositions include
\nolinkurl{WeightedMaximumTarget}, \nolinkurl{ResidualSeparationTarget},
\nolinkurl{ProducerTarget}, \nolinkurl{NeumannCertificateTarget},
\nolinkurl{AutomaticWeightsTarget}, \nolinkurl{PosteriorErrorTarget},
\nolinkurl{DelayedRunTarget} and \nolinkurl{AsynchronousProducerTarget}.
They expose all assumptions. For example, the producer target is
the conjunction of box membership, the finite residual estimate,
and regularity of the explicitly defined sampler. It does not
conclude that some unspecified function exists.

These host functions are specifications for the coded backend, not
an assertion that every Lean function is PR. The object backend
replaces \texttt{Fin d -> Rat} by a length-\(d\) sequence code,
and each evaluator parameter by a registered PR body with its
representation equations. The definitions in \cref{sec:arithmetic}
specify the full numerical representation needed for that step.

\clearpage
\subsection{Core theorem dependency order}

\begin{longtable}{@{}>{\raggedright\arraybackslash}p{0.21\textwidth}>{\raggedright\arraybackslash}p{0.37\textwidth}>{\raggedright\arraybackslash}p{0.34\textwidth}@{}}
\toprule
Target & Definitions and prerequisite lemmas & Proof to implement\\
\midrule
\endhead
Weighted comparison & Rational order; finite maximum with attaining index;
nonnegative matrix multiplication & The positive-maximum contradiction in
\cref{lem:weighted-maximum}, after denominator clearing\\
Residual separation & Weighted triangle inequality; mixed action contract &
Three intermediate vectors and cancellation of \((s-1)d\)\\
Geometric clock & Natural powers; Bernoulli bound; rational division &
QF induction on \(q^n(1+n/s)\le1\), then on block count\\
Regular producer & PR orbit; box bit; diameter; affine envelope &
QF orbit invariant followed by substitution of \(P(k),N(k)\)\\
Neumann enclosure & Matrix powers; finite sum identity &
QF induction on the prefix identity, then weighted comparison\\
Inverse sampler & Induced matrix norm; submultiplicativity; powers &
Finite two-sided telescoping and geometric tail\\
Polynomial chart & Sparse polynomial evaluation, substitution and
formal derivative & Monomial factorization, exact coefficient cancellations,
then the explicit radius inequalities\\
Delayed execution & Finite history; read-time table; update mask &
QF history-window invariant and the bounded coverage bit\\
Refinement coherence & Coded restrictions; nonexpansiveness; residuals &
One rational cancellation, then finite telescoping over model levels\\
Lattice sampler & Dyadic floor; uniform polynomial evaluator; prefix &
Box, contraction, commutation and rounding lemmas, followed by substitution\\
\bottomrule
\end{longtable}

\subsection{The actual induction formulas}

Fix all input parameters and one finite horizon. The following are
the essential invariants, written with rational abbreviations:
\begin{align*}
 \mathcal I_{\mathrm{orbit}}(t):\quad&
 x_t\in D\ \wedge\ x_{t+1}\in D\ \wedge\
 d_w(x_t,x_{t+1})\le q^tR+2s2^{-p};\\
 \mathcal I_{\mathrm{matrix}}(t):\quad&
 A^tw\le q^tw\ \wedge\
 A\sum_{j<t}A^jr+r=\sum_{j<t}A^jr+A^tr;\\
 \mathcal I_{\mathrm{round}}(t):\quad&
 v_t,z_t\in D_n\ \wedge\ d_n(v_t,z_t)\le2h;\\
 \mathcal I_{\mathrm{tower}}(t):\quad&
 d_n(y_n,P_{n,n+t}y_{n+t})\le
       \sum_{j=n}^{n+t-1}e_j.
\end{align*}
Every coordinate conjunction is bounded by a dimension, and every
sum is a PR term. Replace the bounded predicates by their
characteristic-bit equations and clear positive denominators.
The resulting formulas are quantifier-free. If hypotheses are
supplied only through a finite trace, conjoin the corresponding
prefix-validation bit or use its implication form (6).

For a delayed block, the state includes the finite history, the
block counter, and a Boolean mask for coordinates already updated
in that block. The invariant says that all entries in the last
delay window satisfy the old bound, that updated coordinates
satisfy the improved bound, and that all states remain in the
Cartesian box. All time quantifiers range over the finite stored
history. The coverage hypothesis is another bounded bit.

For the tower, compute \(P_{n,m}(z)\) by a loop of length
\(m-n\) that applies \(P_{m-1},P_{m-2},\ldots,P_n\) in
that order. Its level counter decreases by truncated subtraction,
and its vector code is the current state. This is an ordinary
first-order PR implementation of restriction composition; no
higher-type recursion on arbitrary maps is required.

\subsection{Clearing rational denominators correctly}

The compiled integer envelope theorem is not applied to rational
errors without an encoding step. There are two convenient choices.
One can carry numerator and denominator fields through the proof
and use the rational comparison bit directly. Alternatively, for
one finite horizon \(T\), compute a common positive multiple of
all denominators appearing in the data and the finite trace. A
product is sufficient; a least common multiple is unnecessary.
More precisely, let \(e_t\ge0\) be the rational error for
\(t\le T\), and let \(f_t\ge0\) be the rational forcing term
for \(t<T\). Let \(L_T>0\) be that denominator product, including
the denominator of every \(e_t\) through \(T\) and every
\(f_t\) before \(T\). Then
\[
 a_t=L_Te_t\quad(t\le T),\qquad
 r_t=L_Tf_t\quad(t<T)
\]
are nonnegative integers. In the coded implementation they are PR
integer terms, obtained by integer multiplication and division;
divisibility by each original denominator follows from the product
definition. The terms retain \(T\) and the original input as
numerical parameters. Multiplication by the positive \(L_T\)
preserves each recurrence inequality on this finite horizon.

The guarded-loop interface needs these inequalities only for
\(t<T\). To instantiate the compiled host theorem in its
all-natural-index form, one may instead extend \(a_t\) by zero
for \(t>T\), and \(r_t\) by zero for \(t\ge T\).
The transition at \(T\) then has zero left side, and every later
transition is \(0\le0\). This explicit totalization preserves
the value at \(T\) and requires no common denominator for an
infinite sequence.

For the scalar contraction route, the coefficient is exactly
\((s-1)/s\), so the cleared recurrence has the form
\(s a_{t+1}\le(s-1)a_t+r_t\). This matches the compiled clock
and the constant-forcing lemmas. If rational matrix entries are
cleared directly, the resulting numerator and denominator need
not differ by one. Use the general
\texttt{coupled\_envelope} with arbitrary \(c,d\), or first
pass to the rational weighted maximum and then use the scalar
route. This distinction avoids imposing an accidental integrality
restriction on the matrix theorem.

The denominator product may be large, but is a PR function of the
finite trace. The proof does not require a complexity bound from
a stronger arithmetic theory. In the dyadic lattice implementation,
the working denominators are particularly simple because each
state is rounded back to its declared dyadic grid.

\subsection{Object-level packaging and discharged assumptions}

The adapter's \texttt{Package} contains actual PR terms for initial
state, step and fuel; a bounded invariant syntax; a guard and
result formula; and a \texttt{GuardedLoop.Loop} certificate.
The certificate requires actual object proofs of the base,
preservation, and final implication. The export has the form
\begin{lstlisting}
Gamma |- [PRAAx] guard(u)
------------------------------------------
Gamma |- [PRAAx] result(Run(init, step, fuel(u), u), u)
\end{lstlisting}
The displayed notation is explanatory; the file uses the repository's
literal judgment and constructors.

For the general approximate-action theorem, the box and mixed-action
contracts remain assumptions in the sequent. For the polynomial
frontend they are discharged by coefficient algebra. For the
infinite lattice, the concrete polynomial, weight, floor and
boundary estimates discharge every such analytic assumption.
The remaining premises in an object proof are then only its
ordinary numerical input conditions, such as positive dimension.
The construction should export the actual universally quantified
rational-code sentences listed in \cref{sec:lattice}.

No global principle ``every accepted proof trace is true'' is
postulated. The existing \texttt{Prf.trans}, substitution and
loop theorems compose proof values supplied at the relevant
formulas. The elementary arithmetic obligations still need to
be encoded and checked in that calculus.

\subsection{Semantic interpretation is an additional theorem}

The principal iPRA outputs are finite rational estimates. In an
ordinary semantic layer one may interpret their regular samplers
as real vectors, coordinate functions or inverse-limit objects.
The norm estimates and residuals prove that these interpreted
objects satisfy the advertised equations. That interpretation
does not strengthen the object induction rule, nor does a proof
about external real numbers automatically return an object
derivation.

A complete Lean development should therefore expose two named
endpoints for each application: its literal coded iPRA sentence
with a \texttt{Prf PRAAx} proof, and its semantic adequacy theorem
in the chosen represented-object library. The present packet
contains the human proofs and the detailed construction path;
it does not claim those two final layers have already been
implemented in full.

\section{Reproduction, numerical evidence, and remaining work}
\label{sec:validation}

\subsection{What was actually run}

The arithmetic module was compiled with Lean 4.32.0, and
\texttt{\#print axioms} was run for all eighteen theorem declarations.
The recorded dependencies are subsets of \texttt{propext} and
\texttt{Quot.sound}. There is no \texttt{sorryAx} in those reports.
The rational blueprint definitions were also typechecked; their
named theorem targets are definitions of propositions, not proved
theorems. The repository adapter was source-inspected against the
pinned APIs and is explicitly marked uncompiled.

The main Python reference suite ran twenty-four unittest methods
using only integers and \texttt{fractions.Fraction}. It covers
valid and invalid gain certificates, positive-weight construction,
finite Neumann supersolutions, the affine example, floor rounding,
the nonlinear tower through dimension sixteen, multilevel
restriction, same-precision equation residuals, and the alternating
fixed-point counterexample. Boundary tests include zero dimension,
the contraction endpoint \(s=1\), empty Neumann prefixes,
malformed arrays and rejection of floating-point certificate data.

A second reference script checks the approximate-action and
delayed-execution examples, including 3,564 exact action-contract
instances. A third script verifies 1,200 regular-name comparisons
and 120 equation-residual instances for the upper branch, using
only rational bisection and backward substitution. All these finite
computations are independent checks on
constants, indexing, rounding and the implementation. They are
not substituted for the universal human proofs or for object-level
iPRA derivations.

\subsection{Sample coordinates with certified enclosure radii}

The following entries are exact outputs of (L8) at \(k=10\).
The universal proof of (L9), interpreted as a name, places the
represented coordinate within \(2^{-12}=1/4096\) of each sample.
The computations verify the finite samples and their residuals;
the enclosure interpretation uses the theorem.

\begin{center}
\begin{tabular}{crr}
\toprule
Coordinate & Exact rational sample & Decimal sample\\
\midrule
0 & \(580853/2097152\) & \(0.276972293853759765625\)\\
1 & \(145821/1048576\) & \(0.13906574249267578125\)\\
2 & \(48857/524288\) & \(0.0931873321533203125\)\\
3 & \(1176899/16777216\) & \(0.070148646831512451171875\)\\
\bottomrule
\end{tabular}
\end{center}

The CSV and JSON files also give exact interval endpoints, model
dimensions, iteration counts, working precisions and checked
finite residuals. No decimal rounding participates in the
inequality checks.

\subsection{Rebuild instructions}

From the packet root, run
\begin{lstlisting}
python3 code/run_checks.py
python3 code/approximate_actions.py
python3 code/upper_branch_reference.py \
  --output validation/upper_branch_checks.json
bash lean/build.sh /absolute/path/to/lean
bash build_pdf.sh
\end{lstlisting}
Python needs only its standard library. The Lean checks require
the supplied version in \texttt{lean-toolchain}; the large
toolchain distribution is not included. The PDF uses standard
LaTeX packages and is built from the supplied complete source.
The root README describes the separate ProveIt integration
requirements for the uncompiled object adapter.

\subsection{Priority order for repository integration}

The next substantive integration is the rational and finite-vector
frontend for the existing ordinary arithmetic library, followed
by the weighted residual theorem and one complete lattice
endpoint. That sequence yields a concrete theorem of the target
object theory before adding additional applications. The delayed
solver and matrix inverse can then share the same rational
representation and matrix proofs.

The chart producer supplies a useful subsequent application in
multivariate local algebra. Extending the tower theorem to a
particular spatial, spectral or quadrature method requires an
actual commutation estimate and a summable defect-to-gap budget
for that method. The framework identifies those obligations
precisely; it does not turn their absence into a proof.

The outstanding effective height, analytic-number-theory and
singular-geometry obligations discussed in the source audit remain
outstanding. The completed contribution is a reusable mathematical
producer for a substantial numerical feedback and refinement
subclass, with explicit algorithms, complete finite proofs, and a
verified arithmetic foundation for the requested Lean development.


\clearpage
\begin{thebibliography}{99}
\raggedright
\bibitem{audit-unification} ProveIt repository, \emph{ProveIt iPRA unification plan}.
Pinned commit \texttt{3ee850b4c2d650cf01263823fd5415963747c5d0}.
\url{https://github.com/VladimirReshetnikov/ProveIt/blob/3ee850b4c2d650cf01263823fd5415963747c5d0/Logic/PrimitiveRecursiveArithmetic/UNIFICATION-PLAN.md}.

\bibitem{audit-progress} ProveIt repository, \emph{Current iPRA derivation progress inventory}.
Pinned commit \texttt{3ee850b4c2d650cf01263823fd5415963747c5d0}.
\url{https://github.com/VladimirReshetnikov/ProveIt/blob/3ee850b4c2d650cf01263823fd5415963747c5d0/Logic/PrimitiveRecursiveArithmetic/Derivations/PROGRESS.md}.

\bibitem{audit-scalar} ProveIt repository, \emph{A positive repair for Puiseux data: a separated simple branch}.
Pinned commit \texttt{3ee850b4c2d650cf01263823fd5415963747c5d0}.
\url{https://github.com/VladimirReshetnikov/ProveIt/blob/3ee850b4c2d650cf01263823fd5415963747c5d0/Logic/PrimitiveRecursiveArithmetic/docs/proofs/foundations/ipra_embedding_framework/paper/03_analysis.tex}.

\bibitem{audit-scalar-producer} ProveIt repository, \emph{Producing a certificate from a supplied simple-root chart}.
Pinned commit \texttt{3ee850b4c2d650cf01263823fd5415963747c5d0}.
\url{https://github.com/VladimirReshetnikov/ProveIt/blob/3ee850b4c2d650cf01263823fd5415963747c5d0/Logic/PrimitiveRecursiveArithmetic/docs/proofs/foundations/ipra_framework/paper/sections/06b_chart_producer.tex}.

\bibitem{audit-ode} ProveIt repository, \emph{Polynomial differential equations from finite arithmetic}.
Pinned commit \texttt{3ee850b4c2d650cf01263823fd5415963747c5d0}.
\url{https://github.com/VladimirReshetnikov/ProveIt/blob/3ee850b4c2d650cf01263823fd5415963747c5d0/Logic/PrimitiveRecursiveArithmetic/docs/proofs/foundations/ipra_local_global/paper/04_ode.tex}.

\bibitem{audit-demand} ProveIt repository, \emph{A concrete cyclic two-coordinate demand closure}.
Pinned commit \texttt{3ee850b4c2d650cf01263823fd5415963747c5d0}.
\url{https://github.com/VladimirReshetnikov/ProveIt/blob/3ee850b4c2d650cf01263823fd5415963747c5d0/Logic/PrimitiveRecursiveArithmetic/Derivations/Lean/IPRADerivations/Foundations/TwoCoordinateDemandClosure.lean}.

\bibitem{audit-adaptive} ProveIt repository, \emph{Finite-candidate state-adaptive interface calls}.
Pinned commit \texttt{3ee850b4c2d650cf01263823fd5415963747c5d0}.
\url{https://github.com/VladimirReshetnikov/ProveIt/blob/3ee850b4c2d650cf01263823fd5415963747c5d0/Logic/PrimitiveRecursiveArithmetic/Derivations/Lean/IPRADerivations/Foundations/RealizerStateAdaptive.lean}.

\bibitem{audit-fiber} ProveIt repository, \emph{FIBER: effective analysis and adjacent-motion bounds}.
Pinned commit \texttt{3ee850b4c2d650cf01263823fd5415963747c5d0}.
\url{https://github.com/VladimirReshetnikov/ProveIt/blob/3ee850b4c2d650cf01263823fd5415963747c5d0/Logic/PrimitiveRecursiveArithmetic/docs/proofs/foundations/fiber_ipra_2026-11-16/paper/sections/07_analysis.tex}.

\bibitem{audit-coherent} ProveIt repository, \emph{Coherent finite presentations}.
Pinned commit \texttt{3ee850b4c2d650cf01263823fd5415963747c5d0}.
\url{https://github.com/VladimirReshetnikov/ProveIt/blob/3ee850b4c2d650cf01263823fd5415963747c5d0/Logic/PrimitiveRecursiveArithmetic/docs/proofs/foundations/ipra_coherent_methods_2026-11-15/paper/sections/03_refinement.tex}.

\bibitem{audit-topology} ProveIt repository, \emph{Certified observation towers}.
Pinned commit \texttt{3ee850b4c2d650cf01263823fd5415963747c5d0}.
\url{https://github.com/VladimirReshetnikov/ProveIt/blob/3ee850b4c2d650cf01263823fd5415963747c5d0/Logic/PrimitiveRecursiveArithmetic/docs/proofs/topology/ipra_topology_framework_2026-11-15/paper/sections/04_observation_towers.tex}.

\bibitem{audit-finer} ProveIt repository, \emph{Stable observers and finite observational descent}.
Pinned commit \texttt{3ee850b4c2d650cf01263823fd5415963747c5d0}.
\url{https://github.com/VladimirReshetnikov/ProveIt/blob/3ee850b4c2d650cf01263823fd5415963747c5d0/Logic/PrimitiveRecursiveArithmetic/docs/proofs/foundations/FINER-iPRA-framework-2026-11-16/paper/sections/09_observational_descent.tex}.

\bibitem{audit-faltings} ProveIt repository, \emph{Faltings' theorem and intuitionistic PRA: scope and status}.
Pinned commit \texttt{3ee850b4c2d650cf01263823fd5415963747c5d0}.
\url{https://github.com/VladimirReshetnikov/ProveIt/blob/3ee850b4c2d650cf01263823fd5415963747c5d0/Logic/PrimitiveRecursiveArithmetic/docs/proofs/diophantine-number-theory/faltings_ipra/README.md}.

\bibitem{audit-siegel} ProveIt repository, \emph{Siegel: the general theorem and a precise remaining route}.
Pinned commit \texttt{3ee850b4c2d650cf01263823fd5415963747c5d0}.
\url{https://github.com/VladimirReshetnikov/ProveIt/blob/3ee850b4c2d650cf01263823fd5415963747c5d0/Logic/PrimitiveRecursiveArithmetic/docs/proofs/diophantine-number-theory/siegel_ipra/paper/05_general_route.tex}.

\bibitem{audit-primes} ProveIt repository, \emph{Siegel--Walfisz and Bombieri--Vinogradov targets and finite ingredients}.
Pinned commit \texttt{3ee850b4c2d650cf01263823fd5415963747c5d0}.
\url{https://github.com/VladimirReshetnikov/ProveIt/blob/3ee850b4c2d650cf01263823fd5415963747c5d0/Logic/PrimitiveRecursiveArithmetic/docs/proofs/analytic-number-theory/prime_distribution_ipra/paper/sections/07_sw_bv.tex}.


\bibitem{rump1994} S. M. Rump, \emph{Verification methods for dense and
sparse systems of equations}, in J. Herzberger (ed.), Topics in Validated
Computations, Elsevier, 1994, pp. 63--136.
\url{https://www.tuhh.de/ti3/paper/rump/Ru94.pdf}.
\bibitem{radii} O. H\'enot and contributors, \emph{The radii polynomial
approach}, RadiiPolynomial.jl documentation, accessed 19 November 2026.
\url{https://olivierhnt.github.io/RadiiPolynomial.jl/stable/radii_polynomial_approach/}.
\bibitem{fefermanstrahm} S. Feferman and T. Strahm,
\emph{Unfolding finitist arithmetic}, manuscript dated 1 January 2010.
\url{https://math.stanford.edu/~feferman/papers/UnfoldFA.pdf}.
\bibitem{kohlenbachsafarik} U. Kohlenbach and P. Safarik,
\emph{Fluctuations, effective learnability and metastability in analysis}.
\url{https://www2.mathematik.tu-darmstadt.de/~kohlenbach/mind_ps_revised.pdf}.
\bibitem{leanrelease} Lean project, \emph{Lean 4.32.0 release notes}.
\url{https://lean-lang.org/doc/reference/latest/releases/v4.32.0/}.
\end{thebibliography}
\end{document}
